%% Copyright 2012-<COPYRIGHT_YEAR> by abnTeX2 group at http://www.abntex.net.br/ 
%%
%% This work may be distributed and/or modified under the
%% conditions of the LaTeX Project Public License, either version 1.3
%% of this license or (at your option) any later version.
%% The latest version of this license is in
%%   http://www.latex-project.org/lppl.txt
%% and version 1.3 or later is part of all distributions of LaTeX
%% version 2005/12/01 or later.
%%
%% This work has the LPPL maintenance status `maintained'.
%% 
%% The Current Maintainer of this work is the abnTeX2 team, led
%% by Lauro César Araujo. Further information are available on 
%% http://www.abntex.net.br/
%%
%% This work consists of the files abntex2-modelo-trabalho-academico.tex,
%% abntex2-modelo-include-comandos and abntex2-modelo-references.bib
%%

% ------------------------------------------------------------------------
% ------------------------------------------------------------------------
% abnTeX2: Modelo de Trabalho Academico (tese de doutorado, dissertacao de
% mestrado e trabalhos monograficos em geral) em conformidade com 
% ABNT NBR 14724:2011: Informacao e documentacao - Trabalhos academicos -
% Apresentacao
% ------------------------------------------------------------------------
% ------------------------------------------------------------------------

\documentclass[
	% -- opções da classe memoir --
	12pt,				% tamanho da fonte
	openright,			% capítulos começam em pág ímpar (insere página vazia caso preciso)
	twoside,			% para impressão em recto e verso. Oposto a oneside
	a4paper,			% tamanho do papel. 
	% -- opções da classe abntex2 --
	%chapter=TITLE,		% títulos de capítulos convertidos em letras maiúsculas
	%section=TITLE,		% títulos de seções convertidos em letras maiúsculas
	%subsection=TITLE,	% títulos de subseções convertidos em letras maiúsculas
	%subsubsection=TITLE,% títulos de subsubseções convertidos em letras maiúsculas
	% -- opções do pacote babel --
	english,			% idioma adicional para hifenização
	french,				% idioma adicional para hifenização
	spanish,			% idioma adicional para hifenização
  japanese,
	brazil				% o último idioma é o principal do documento
	]{abntex2}

% ---
% Pacotes 
% ---
\usepackage{lmodern}			% Usa a fonte Latin Modern			
\usepackage[T1]{fontenc}		% Selecao de codigos de fonte.
\usepackage[utf8]{inputenc}		% Codificacao do documento (conversão automática dos acentos)
\usepackage{indentfirst}		% Indenta o primeiro parágrafo de cada seção.
\usepackage{color}				% Controle das cores
\usepackage{graphicx}			% Inclusão de gráficos
\usepackage{microtype} 			% para melhorias de justificação
\usepackage{pdfpages}           % necessário para comando \includepdf
\usepackage{lipsum}				% para geração de dummy text
\usepackage{amsmath}
\usepackage[braket, qm]{qcircuit}
\usepackage[version=4]{mhchem}
\usepackage{multirow}

% Force titles to use Latin Modern (serif)
\renewcommand{\ABNTEXchapterfont}{\normalfont\rmfamily\bfseries}

% ---

% ---
% Pacotes de citações
% ---
% Opção 1 -> Mais bonito e limpo, mas na sessão de referências, não fica da forma correta;
% \usepackage[brazilian,hyperpageref]{backref}
% \usepackage{cite}
% \usepackage{hyperref}
% \bibliographystyle{plain}

% Opção 2 -> Provavelmente é a mais correta. Melhor deixar aqui caso eu precise mudar rapidamente para ficar dentro das regras;
\usepackage[brazilian,hyperpageref]{backref}	 % Paginas com as citações na bibl
\usepackage[num]{abntex2cite}	% Citações padrão ABNT
\citebrackets[]

% --- 
% CONFIGURAÇÕES DE PACOTES
% --- 

% ---
% Configurações do pacote backref
% Usado sem a opção hyperpageref de backref
\renewcommand{\backrefpagesname}{Citado na(s) página(s):~}
% Texto padrão antes do número das páginas
\renewcommand{\backref}{}
% Define os textos da citação
\renewcommand*{\backrefalt}[4]{
	\ifcase #1 %
		Nenhuma citação no texto.%
	\or
		Citado na página #2.%
	\else
		Citado #1 vezes nas páginas #2.%
	\fi}%
% ---

% ---
% Variáveis utilizadas ao longo do texto
% ---
\newcommand{\aprov}{20 de junho de 2025}

% ---
% Informações de dados para CAPA e FOLHA DE ROSTO
% ---
\titulo{Análise comparativa de modelos de Predição do Nível da Água no Rio Guaíba}
\autor{Juliano Almeida Machado}
\local{Porto Alegre, Brasil}
\data{2025}
\orientador{Sebastián Gonçalves}
% \coorientador{Equipe \abnTeX}
\instituicao{%
  Universidade Federal do Rio Grande do Sul -- UFRGS
  \par
  Instituto de Física}
\tipotrabalho{Trabalho (Conclusão de Curso)}
% O preambulo deve conter o tipo do trabalho, o objetivo, 
% o nome da instituição e a área de concentração 
% \preambulo{Modelo canônico de trabalho monográfico acadêmico em conformidade com
% as normas ABNT apresentado à comunidade de usuários}
\preambulo{Trabalho de Conclusão apresentado à Comissão de Graduação dos Cursos de Física da Universidade Federal do Rio
 Grande do Sul, como requisito parcial para obtenção do título de Bacharel em Física.}
% ---


% ---
% Configurações de aparência do PDF final

% alterando o aspecto da cor azul
\definecolor{blue}{RGB}{41,5,195}

% informações do PDF
\makeatletter
\hypersetup{
     	%pagebackref=true,
		pdftitle={\@title}, 
		pdfauthor={\@author},
    	pdfsubject={\imprimirpreambulo},
	    pdfcreator={LaTeX with abnTeX2},
		pdfkeywords={abnt}{latex}{abntex}{abntex2}{trabalho acadêmico}, 
		colorlinks=true,       		% false: boxed links; true: colored links
    	linkcolor=blue,          	% color of internal links
    	citecolor=blue,        		% color of links to bibliography
    	filecolor=magenta,      		% color of file links
		urlcolor=blue,
		bookmarksdepth=4
}
\makeatother
% --- 

% ---
% Posiciona figuras e tabelas no topo da página quando adicionadas sozinhas
% em um página em branco. Ver https://github.com/abntex/abntex2/issues/170
\makeatletter
\setlength{\@fptop}{5pt} % Set distance from top of page to first float
\makeatother
% ---

% ---
% Possibilita criação de Quadros e Lista de quadros.
% Ver https://github.com/abntex/abntex2/issues/176
%
\newcommand{\quadroname}{Quadro}
\newcommand{\listofquadrosname}{Lista de quadros}

\newfloat[chapter]{quadro}{loq}{\quadroname}
\newlistof{listofquadros}{loq}{\listofquadrosname}
\newlistentry{quadro}{loq}{0}

% configurações para atender às regras da ABNT
\setfloatadjustment{quadro}{\centering}
\counterwithout{quadro}{chapter}
\renewcommand{\cftquadroname}{\quadroname\space} 
\renewcommand*{\cftquadroaftersnum}{\hfill--\hfill}

\setfloatlocations{quadro}{hbtp} % Ver https://github.com/abntex/abntex2/issues/176
% ---

% --- 
% Espaçamentos entre linhas e parágrafos 
% --- 

% O tamanho do parágrafo é dado por:
\setlength{\parindent}{1.3cm}

% Controle do espaçamento entre um parágrafo e outro:
\setlength{\parskip}{0.2cm}  % tente também \onelineskip

% ---
% compila o indice
% ---
\makeindex
% ---

% ----
% Início do documento
% ----
\begin{document}
\renewcommand{\familydefault}{\rmdefault} 

% Seleciona o idioma do documento (conforme pacotes do babel)
%\selectlanguage{english}
\selectlanguage{brazil}

% Retira espaço extra obsoleto entre as frases.
\frenchspacing 

% ----------------------------------------------------------
% ELEMENTOS PRÉ-TEXTUAIS
% ----------------------------------------------------------
% \pretextual

% ---
% Capa
% ---
\imprimircapa
% ---

% ---
% Folha de rosto
% (o * indica que haverá a ficha bibliográfica)
% ---
\imprimirfolhaderosto*
% ---

% ---
% Inserir a ficha bibliografica
% ---

% Isto é um exemplo de Ficha Catalográfica, ou ``Dados internacionais de
% catalogação-na-publicação''. Você pode utilizar este modelo como referência. 
% Porém, provavelmente a biblioteca da sua universidade lhe fornecerá um PDF
% com a ficha catalográfica definitiva após a defesa do trabalho. Quando estiver
% com o documento, salve-o como PDF no diretório do seu projeto e substitua todo
% o conteúdo de implementação deste arquivo pelo comando abaixo:
%
% \begin{fichacatalografica}
%     \includepdf{fig_ficha_catalografica.pdf}
% \end{fichacatalografica}

\begin{fichacatalografica}
	\sffamily
	\vspace*{\fill}					% Posição vertical
	\begin{center}					% Minipage Centralizado
	\fbox{\begin{minipage}[c][8cm]{15.5cm}		% Largura -> Tive que aumentar o tamanho da caixa. Original 13.5cm;
	\small
	\imprimirautor
	%Sobrenome, Nome do autor
	
	\hspace{0.5cm} \imprimirtitulo  / \imprimirautor. --
	\imprimirlocal, \imprimirdata-
	
	\hspace{0.5cm} \thelastpage p. : il. (algumas color.) ; 30 cm.\\
	
	\hspace{0.5cm} \imprimirorientadorRotulo~\imprimirorientador\\
	
	\hspace{0.5cm}
	\parbox[t]{\textwidth}{\imprimirtipotrabalho~--~\imprimirinstituicao,
	\imprimirdata.}\\
	
	\hspace{0.5cm}
		1. Séries Temporais.
		2. Enchentes.
		3. Guaíba.
		I. Orientador.
		II. Universidade Federal do Rio Grande do Sul.
		III. Instituto de Física.
		IV. Título 			
	\end{minipage}}
	\end{center}
\end{fichacatalografica}

% ---
% Inserir folha de aprovação
% ---

% Isto é um exemplo de Folha de aprovação, elemento obrigatório da NBR
% 14724/2011 (seção 4.2.1.3). Você pode utilizar este modelo até a aprovação
% do trabalho. Após isso, substitua todo o conteúdo deste arquivo por uma
% imagem da página assinada pela banca com o comando abaixo:
%
% \begin{folhadeaprovacao}
% \includepdf{folhadeaprovacao_final.pdf}
% \end{folhadeaprovacao}
%
% \begin{folhadeaprovacao}

%   \begin{center}
%     {\ABNTEXchapterfont\large\imprimirautor}

%     \vspace*{\fill}\vspace*{\fill}
%     \begin{center}
%       \ABNTEXchapterfont\bfseries\Large\imprimirtitulo
%     \end{center}
%     \vspace*{\fill}
    
%     \hspace{.45\textwidth}
%     \begin{minipage}{.5\textwidth}
%         \imprimirpreambulo
%     \end{minipage}%
%     \vspace*{\fill}
%    \end{center}
        
%    Trabalho aprovado. \imprimirlocal, \aprov:

%    \assinatura{\textbf{\imprimirorientador} \\ Orientador} 
%    \assinatura{\textbf{Professor} \\ Convidado 1}
%    \assinatura{\textbf{Professor} \\ Convidado 2}
%    %\assinatura{\textbf{Professor} \\ Convidado 3}
%    %\assinatura{\textbf{Professor} \\ Convidado 4}
      
%    \begin{center}
%     \vspace*{0.5cm}
%     {\large\imprimirlocal}
%     \par
%     {\large\imprimirdata}
%     \vspace*{1cm}
%   \end{center}
  
% \end{folhadeaprovacao}
% ---

% ---
% Agradecimentos
% ---
% \begin{agradecimentos}
% \lipsum[1]

% \end{agradecimentos}
% ---

% ---
% RESUMOS
% ---

% resumo em português
% \setlength{\absparsep}{18pt} % ajusta o espaçamento dos parágrafos do resumo
% \begin{resumo}
% Desastres naturais associados a eventos climáticos extremos têm impactado severamente regiões vulneráveis, como
% evidenciado pelas enchentes que ocorreram no estado brasileiro do Rio Grande do Sul em 2024. Este trabalho propõe uma
% análise comparativa de modelos de predição do nível da água do Rio Guaíba, visando aprimorar sistemas de alerta precoce
% para gestão de crises. Foram avaliados quatro modelos baseados em dados: SARIMA (estatístico), LSTM (rede neural
% recorrente), XGBoost e LightGBM (baseados em Gradient Boosting), utilizando séries temporais horárias de nível d'água e
% variáveis meteorológicas (2019-2024) do Sistema Nacional de Informações sobre Recursos Hídricos (SNIRH). Os dados foram
% pré-processados com interpolação de falhas e normalização, e os modelos otimizados via validação cruzada temporal. O
% desempenho foi quantificado por métricas hidrologicamente relevantes: RMSE (raiz do erro quadrático médio), MAE (erro
% absoluto médio) e NSE (coeficiente de eficiência de Nash-Sutcliffe). 

% % Resultados preliminares indicam que o LightGBM obteve o melhor equilíbrio entre acurácia (RMSE = 8,3 cm; NSE = 0,96) e
% % eficiência computacional (< 2 segundos para previsões em tempo real), superando modelos tradicionais em cenários de
% % variação abrupta. Conclui-se que modelos baseados em boosting são alternativas viáveis para sistemas de alerta
% % integrados, reduzindo a dependência de simulações hidrodinâmicas complexas e democratizando o acesso a informações
% % críticas para tomada de decisão em eventos extremos.

%  \textbf{Palavras-chave}: Predição de enchentes; Séries temporais; Rio Guaíba; Machine Learning; Gestão de desastres;
% \end{resumo}

% % resumo em inglês
% \begin{resumo}[Abstract]
%  \begin{otherlanguage*}{english}

% Natural disasters linked to extreme weather events have severely impacted vulnerable regions, notably the 2024 floods in
% Rio Grande do Sul, Brazil. This study presents a comparative analysis of water-level prediction models for the Guaíba
% River, aiming to enhance early-warning systems for crisis management. Four data-driven models were evaluated: SARIMA
% (statistical), LSTM (recurrent neural network), XGBoost, and LightGBM (gradient-boosting-based), using hourly
% water-level time series and meteorological variables (2019-2024) from the National Water Resources Information System
% (SNIRH). Data underwent gap-filling interpolation and normalization, with models optimized via temporal
% cross-validation. Performance was assessed using hydrologically relevant metrics: RMSE (root mean squared error), MAE
% (mean absolute error), and NSE (Nash-Sutcliffe efficiency coefficient).

% % Preliminary results indicate that LightGBM achieved the best balance between accuracy (RMSE = 8.3 cm; NSE = 0.96) and
% % computational efficiency (< 2 seconds for real-time predictions), outperforming traditional models in scenarios of
% % abrupt variation. We conclude that boosting-based models are viable alternatives for integrated alert systems, reducing
% % reliance on complex hydrodynamic simulations and democratizing access to critical information for decision-making during
% % extreme events.

%    \vspace{\onelineskip}
 
%    \noindent 
%    \textbf{Keywords}: Flood prediction; Time series; Guaíba River; Machine Learning; Disaster management;
%  \end{otherlanguage*}
% \end{resumo}

% % ---
% % inserir lista de ilustrações
% % ---
% \pdfbookmark[0]{\listfigurename}{lof}
% \listoffigures*
% \cleardoublepage
% % ---

% % ---
% % inserir lista de quadros
% % ---
% \pdfbookmark[0]{\listofquadrosname}{loq}
% \listofquadros*
% \cleardoublepage
% % ---

% % ---
% % inserir lista de tabelas
% % ---
% \pdfbookmark[0]{\listtablename}{lot}
% \listoftables*
% \cleardoublepage
% % ---

% ---
% inserir lista de abreviaturas e siglas
% ---
% (EXPANDIR - Revisar todas as Siglas antes de entregar);
\begin{siglas}
  \item[ABNT] Associação Brasileira de Normas Técnicas
  \item[abnTeX] ABsurdas Normas para TeX
  \item[ANA] Agência Nacional de Águas
  \item[ANFIS] Adaptive Neuro Fuzzy Inference System (Sistema de inferência neuro-borroso adaptativo) 
  \item[ARIMA] AutoRegressive Integrated Moving Average (Média Móvel Integrada Auto-Regressiva)
  \item[CEMADEN] Centro Nacional de Monitoramento e Alertas de Desastres Naturais 
  \item[CNN] Rede Neural Convolucional
  \item[CUDA] Compute Unified Device Architecture (Arquitetura de Dispositivo de Computação Unificada)
  \item[DMAE] Departamento Municipal de Água e Esgotos
  \item[ETA] Estação de Tratamento de Água 
  \item[FFNN] Feed Forward Neural Network (Rede Neural Pré-alimentada)
  \item[GBM] Gradient Boosting Machine (Impulsionamento de Gradiente de Máquina) 
  \item[IBGE] Instituto Brasileiro de Geografia e Estatística
  \item[INMET] Instituto Nacional de Meteorologia
  \item[IPCC] Painel Intergovernamental sobre Mudanças Climáticas
  \item[IPH] Instituto de Pesquisas Hidráulicas
  \item[LSTM] Long Short Term Memory (Memória de Longo Curto Prazo)
  \item[LightGBM] Light Gradient-Boosting Machine (Impulsionamento de Gradiente de Máquina Leve) 
  \item[MLP] Multilayer Perceptron (Perceptron de Múltiplas Camadas)
  \item[RMPOA] Região Metropolitana de Porto Alegre 
  \item[RNN] Rede Neural Recorrente
  \item[SARIMA] Seasonal AutoRegressive Integrated Moving Average (Média Móvel Integrada Auto-Regressiva com Sazonalidade)
  \item[SNIRH] Sistema Nacional de Informações sobre Recursos Hídricos
  \item[SVR] Support Vector Regression (Regressão de Vetores de Suporte)
  \item[UFRGS] Universidade Federal do Rio Grande do Sul
  \item[UNFCCC] United Nations Framework Convention on Climate Change (Convenção-Quadro das Nações Unidas sobre a Mudança do Clima)
  \item[XGBoost] eXtreme Gradient Boosting (Impulsionamento de Gradiente Extremo)

\end{siglas}
% ---

% % ---
% % inserir lista de símbolos
% % ---
% \begin{simbolos}
%   \item[$ \Gamma $] Letra grega Gama
%   \item[$ \Lambda $] Lambda
%   \item[$ \zeta $] Letra grega minúscula zeta
%   \item[$ \in $] Pertence
% \end{simbolos}
% % ---

% ---
% inserir o sumario
% ---
\pdfbookmark[0]{\contentsname}{toc}
\tableofcontents*
\cleardoublepage
% ---

% ----------------------------------------------------------
% ELEMENTOS TEXTUAIS
% ----------------------------------------------------------
\textual

% ----------------------------------------------------------
% Introdução (exemplo de capítulo sem numeração, mas presente no Sumário)
% ----------------------------------------------------------
\chapter{Introdução}
Desastres naturais provocados por eventos climáticos extremos geram impactos profundos. No âmbito econômico, recursos
destinados à recuperação de áreas afetadas enfraquecem a microeconomia local e ampliam a vulnerabilidade futura. No
entanto, o impacto social costuma ser mais severo. Traumas, perdas humanas e rupturas comunitárias deixam danos
duradouros, muitas vezes irreversíveis.

Devido às mudanças climáticas, a frequência e intensidade dos desastres naturais vêm aumentando globalmente, como
resultado, muitas regiões são afetadas por eventos climáticos extremos de maneira acelerada, causando um significativo
despreparo. Entre esses desastres, a ocorrência de enchentes se destacam pelo aumento de extremos de precipitação
globais, como evidenciado por Dunn et al. (2024) \cite{dunn}. As mesmas apresentam uma capacidade de causar destruição
generalizada e afetar áreas densamente povoadas, especialmente em regiões urbanas com infraestrutura inadequada para
gerenciar grandes volumes de água. Em muitos casos, essa inadequação é agravada pela falta de manutenção de sistemas
essenciais, como drenagem urbana ou canais fluviais, tornando as cidades vulneráveis mesmo quando dispõem de estruturas
projetadas para lidar com a situação.

No final de abril e início de maio de 2024, o Rio Grande do Sul registrou precipitações excepcionais que elevaram
drasticamente os níveis das bacias hidrográficas. Na região metropolitana de Porto Alegre (RMPOA), ocorreu a maior cheia
da história do Rio Guaíba, superando o recorde da Cheia de 1941 \cite{dep}. Em 5 de maio, o nível atingiu 5{,}33 metros,
e falhas no sistema de contenção - muros, diques, comportas e estações de bombeamento \cite{comercio} -, resultaram no
alagamento de bairros inteiros, evacuações nas áreas mais afetadas, bloqueio das principais vias e interrupções de
abastecimento, energia e mobilidade. O evento decorreu de fatores meteorológicos combinados à vulnerabilidade urbana
\cite{yang, hammond, zhang, pereira}, deixando milhares de deslocados e centenas de vítimas \cite{alcantara}.

Durante a crise, previsões confiáveis do nível do Guaíba foram essenciais para orientar ações emergenciais. Os modelos
hidrológicos e hidrodinâmicos utilizados, embora precisos, exigem grande volume de dados topográficos, meteorológicos e
hidrológicos, conhecimento especializado e tempo de simulação elevado. Nos momentos críticos, só puderam ser executados
diariamente \cite{iph}, limitando sua utilidade em previsões de curto prazo \cite{atashi}.

Modelos baseados em dados surgem como alternativa mais ágil. Eles utilizam técnicas para identificar e explorar relações
estatísticas entre os dados de entrada e saída. Com a aprendizagem dos padrões presentes nos dados históricos, é
possível mapear as dinâmicas subjacentes do sistema e utilizá-las para prever valores futuros. Assim, esses modelos são
capazes de realizar previsões precisas com base em séries temporais, capturando tendências e comportamentos complexos
que podem ser difíceis de modelar diretamente por meio de abordagens tradicionais. Enquanto esses modelos demandam muito
tempo e dados específicos, os modelos baseados em dados oferecem maior agilidade para previsões em tempo real, sendo uma
escolha mais adequada para sistemas de alerta precoce e mitigação de desastres de cheias \cite{ta}.

Entre os métodos clássicos, destaca-se o ARIMA \cite{box}, que combina autoregressão, média móvel e diferenciação, além
de permitir sazonalidade via SARIMA e variáveis exôgenas através do SARIMAx. Esses modelos têm histórico de aplicação em
hidrologia, como Atashi et al.\ (2022) \cite{atashi} e Güldal et al.\ (2009) \cite{gudal}, servindo como base de
comparação para métodos modernos.

No campo de \textit{Machine Learning}, técnicas de \textit{Gradient Boosting} são amplamente utilizadas pela rapidez e
precisão, como o XGBoost \cite{chen_xgboost} e o LightGBM \cite{lightgbm}, que mantêm desempenho elevado com menor custo
computacional. Em \textit{Deep Learning}, modelos recorrentes como LSTM \cite{hochreiter} capturam dependências
temporais complexas e têm se mostrado eficazes em predições hidrológicas \cite{borwarnginn}.

Este estudo aplica essas abordagens à previsão do nível da água durante as enchentes do Guaíba, avaliando a eficácia de
SARIMA, XGBoost, LSTM e outros métodos em um cenário crítico, onde precisão e rapidez são essenciais. A comparação será
feita por métricas consolidadas, como \textit{Kling–Gupta efficiency} (KGE), \textit{Nash-Sutcliffe Model Efficiency
Coefficient} (NSE), \textit{Root Mean Squared Error} (RMSE) e \textit{Mean Absolute Error} (MAE), permitindo avaliar
tanto a acurácia quanto a capacidade de cada modelo de reproduzir a dinâmica hidrológica.

Também será analisado o uso de variáveis externas, como precipitação, vento e afluentes, além de aplicar transformações
como \textit{lag}, janelas temporais e cumulativos. O objetivo é identificar não apenas o menor erro preditivo, mas o
comportamento de cada técnica frente às diferentes dinâmicas do sistema. Essa análise contribui para estratégias mais
robustas de gestão de desastres.

Para esta análise, serão utilizados dados  obtidos das estações telemétricas relevantes através do Sistema Nacional de
Informações sobre Recursos Hídricos (SNIRH), que disponibiliza informações das principais bacias hidrográficas do Brasil
\cite{snirh}. Esses dados serão a base para as séries históricas do nível das águas do Guaíba, complementados por dados
externos, como precipitação, vazão e temperatura também do SNIRH.

% \part{Referenciais teóricos}
% ---------------------------------------------------------------------------------------------------------------------
% Capitulo de revisão de literatura -> Situação geográfica e metereológica. Modelos de Ciência de Dados. Basicamente serve para adicionar o contexto do problema;
% ---------------------------------------------------------------------------------------------------------------------

\chapter{Contextualização}
Neste capítulo será feita uma análise dos fatores e contexto que levaram à crise das Enchentes de 2024 no estado do Rio
Grande do Sul, com um destaque para o aumento de precipitação causado pelo aquecimento global e seus efeitos,
especialmente na América Latina. Serão exploradas as causas e características de enchentes e os efeitos sociais e
economicos destes eventos, com foco em centros urbanos. Será também explorada a vulnerabilidade intrinseca de centros
urbanos e a necessidade elevada de prevenção e monitoramento para tais situações, sendo realizada através de uma análise
do monitoramento e sistemas de prevenção presentes na Região Metropolitana de Porto Alegre.
    \section{Mudanças climáticas e desastres naturais}
    A ocorrência de enchentes é frequentemente associada a chuvas intensas e ao transbordamento de rios. No entanto, um
fator subjacente e cada vez mais relevante para a intensificação desses eventos é o aquecimento global. O aumento das
temperaturas globais altera o ciclo hidrológico, intensificando a evaporação, modificando os padrões de precipitação e
contribuindo para eventos extremos mais frequentes e intensos. Dessa forma, compreender a relação entre o aquecimento
global e as enchentes é essencial para o desenvolvimento de estratégias eficazes de mitigação e adaptação

\subsection{Aumento de temperatura e emissões globais}
% V. Aumento de temperaturas globais; V. Ação humana no clima; V. Sul Global mais afetado proporcionalmente
% 3. Aumento de heat waves -> Dunn; V. Aumento de precipitação;
% 5. Aumento de Enchentes (precipitação/precipitação acumulada);
% 6. El Nino e La Nina

% Chapter 2.3.1.1.3 (Temperatures during the instrumental period - surface) of the AR6 WGI;
As últimas décadas foram marcadas por quebras constantes de recordes climáticos, especialmente em máximas de
temperatura, evidenciando de forma empírica o impacto das mudanças climáticas em escala global. Esta tendência pode ser
evidenciada através das Anomalias de Temperatura que são definidas como a diferença entre a temperatura observada em um
determinado período e uma temperatura média de referência calculada para um período base. As mesmas são amplamente
utilizadas como indicadores de tendências e mudanças de longo prazo nos padrões climáticos, pois podem capturar
variações consistentes e sistemáticas não aparentes em dados absolutos, permitindo uma análise mais clara das mudanças
globais. Nas últimas décadas, as Anomalias de Temperatura apresentam um padrão de crescimento constante, como pode
evidenciar-se na Figura~\ref{temp_anomaly}.

\begin{figure}[htb]
	\caption{\label{temp_anomaly} Anomalia de Temperatura referente ao período de 1850 a 2024}
	\begin{center}
	    \includegraphics[scale=0.455]{img/Temperature Anomaly.png}
	\end{center}
	\legend{Fonte: \textit{National Oceanic and Atmospheric Administration} (NOAA) \cite{ncei_noaa}}
\end{figure}

De acordo com o sexto - e mais recente - relatório do Painel Intergovernamental sobre Mudanças Climáticas (IPCC),
estabelecido pelas Nações Unidas e amplamente reconhecido como a principal referência na análise do conhecimento
científico relacionado às mudanças climáticas, a superfície terrestre tem sofrido um aumento expressivo de temperatura
em comparação com períodos históricos. Considerando o intervalo de 1850 a 1900 (utilizado como base pré-industrial por
ser o primeiro período com boas observações de temperaturas registradas), a temperatura média da superfície
terrestre\footnote{Dados referentes à medida do GMST - Temperatura Média Global da Superfície - que representa uma média
ponderada entre as temperaturas das massas terrestres e das massas oceânicas.} no período de 2001 a 2020 foi
aproximadamente 1.07ºC mais elevada. Além disso, o aumento da temperatura somente nas últimas quatro décadas (1980 a
2020) é de 0,76ºC, evidenciando uma aceleração no aquecimento global \cite{gulev}.

Esse aumento, aparentemente pequeno, não reflete a real magnitude das modificações no sistema terrestre como um todo. É
importante notar que um aumento de 1ºC representa uma média global; portanto, houveram locais no qual o aquecimento foi
significativamente mais pronunciado, como nos polos. Nessas regiões, mudanças pequenas na temperatura podem levar ao
derretimento de geleiras e, logo, na diminuição do albedo da Terra, isto é, a capacidade de refletir os raios solares de
volta para o espaço. Portanto, nessa situação, a Terra absorverá mais calor, levando à ainda mais derretimento das
geleiras, criando assim um ciclo de retroalimentação. Além disso, é relevante destacar que um aquecimento de 1ºC em toda
a superfície terrestre é extremamente significativo devido à massa do planeta. Para que um aumento dessa magnitude
ocorra em um corpo com a massa da Terra, a quantidade de energia térmica\footnote{De forma bastante simplificada, essa
energia pode ser vista pela fórmula $ Q = mc\Delta T $, onde $m$ é a massa, $c$ é o calor específico, e $\Delta T$ a
variação de temperatura.} liberada é imensa.

% Chapter 3.3 (Human Influence on the Atmosphere and Surface) of the AR6 WGI;
É indubitável que a ação humana foi o maior contribuinte para o Aquecimento Global \cite{ipcc, gillett}. As atividades
industriais que sustentam os padrões de consumo modernos são as principais emissoras de gases responsáveis pelo efeito
estufa, como dióxido de carbono (\ce{CO2}), metano (\ce{CH4}), óxido nitroso (\ce{N2O}), entre outros. Esses gases
possuem a capacidade de absorver a radiação infravermelha emitida pela superfície da Terra, retendo o calor que, de
outra forma, seria dissipado no espaço. 

Embora o efeito estufa seja considerado um fenômeno essencial graças à concentração natural de certos GEEs (sem eles a
temperatura média da superfície terrestre seria significativamente menor), a ação humana tem intensificado esse processo
alarmantemente. Desde o início da Revolução Industrial, a concentração de GEEs na atmosfera aumentou de forma gradual,
principalmente devido à queima de combustíveis fósseis como o carvão, o gás natural e o petróleo. Esses combustíveis são
formados por compostos de carbono de estrutura complexa. Sua alta densidade energética permite que sejam facilmente
queimados para gerar calor, que pode ser convertido em trabalho mecânico em motores, tornando-os a base da matriz
energética global. O principal subproduto desta queima, porém, se dá na forma de gás carbonico (\ce{CO2}) que é liberado
livremente na atmosfera. 

Considerando que a formação de combustíveis fósseis se dá a partir da decomposição de matéria orgânica ao longo de
milhões de anos, a emissão de \ce{CO2} desta forma representa a liberação de carbono na atmosfera que de outra forma
permaneceria preso na superfície. Estima-se que, entre 1850 e 2019, as emissões acumuladas de \ce{CO2} foram de 2400±240
Gt\ce{CO2} \cite{friedlingstein}, resultando em uma concentração de \ce{CO2} na atmosfera de 410 partes por milhão,
comparavél com o início do máximo térmico do Paleoceno-Eoceno \cite{gingerich}, um período associado a uma das mais
recentes extinções em massa.

Esses dados destacam a intensidade das atividades humanas no agravamento do efeito estufa e evidenciam a necessidade
urgente de iniciativas coordenadas em escala global para mitigar os danos climáticos e garantir um futuro mais
sustentável. Iniciativas como o Acordo de Paris são essenciais para garantir a coordenação das ações climáticas ao
estabelecer metas e responsabilidades de cada nação. 

Firmado em 2015 pelas nações participantes do Convenção-Quadro das Nações Unidas sobre a Mudança do Clima (UNFCCC), o
acordo estabelece a necessidade de reduzir as emissões de GEEs em escala global, visando manter o aumento de temperatura
abaixo de 2ºC em relação à base pré-industrial, além de providenciar mecanismos para que nações em desenvolvimento
possam mitigar os efeitos das mudanças climáticas e se adaptar a elas sem comprometer severamente suas atividades
econômicas essenciais.

% Chapter 8.2 of the AR6 WGII and Chapter 12.3.6 of the AR6 WGII;
\subsection{Contribuições e vulnerabilidades climáticas}
\label{vul}
A América Latina é um grande contribuinte na emissão de GEEs, contabilizando aproximadamente 11\% das emissões
históricas globais de \ce{CO2} \cite{friedlingstein}. Contudo, a maior parte dessa emissão vem na forma de uso de terra,
como agricultura, pecuária e desmatamento. Esse padrão é característico de muitas regiões em desenvolvimento, onde a
produção industrial está voltada predominantemente para atender às demandas de exportação de bens e matéria-prima, em
vez de ser direcionada ao consumo interno. Consequentemente, as contribuições dessas regiões para o aquecimento global
vem com o papel de fornecedores globais, beneficiando outras partes do mundo, enquanto enfrentam os desafios associados
a essas emissões.

\begin{figure}[htb]
	\caption{\label{contrib} Emissões Cumulativas por Região (1850-2019)}
	\begin{center}
	    \includegraphics[scale=0.45]{img/contrib_regiao.png}
	\end{center}
	\legend{Fonte: \textit{Global Carbon Budget}, trauduzido pelo autor \cite{friedlingstein}}
\end{figure}

Além disso, essas regiões são significativamente mais vulneráveis aos impactos das mudanças climáticas \cite{chambers},
não apenas por estarem localizadas em áreas tipicamente mais tropicais, onde o aumento das temperaturas é mais intenso,
mas também pela infraestrutura limitada para prevenir e mitigar desastres. Diversos fatores contribuem para a
resiliência de uma região diante das mudanças climáticas, como o desenvolvimento desigual, a fragilidade das
instituições e os altos índices de pobreza. Comunidades mais vulneráveis acabam sendo desproporcionalmente afetadas por
perdas no setor agrícola, o que intensifica a insegurança alimentar e agrava as condições socioeconômicas da região. As
mudanças climáticas também afetam a distribuição dos biomas e sua biodiversidade, criando desafios para as populações
nativas que dependem diretamente desses recursos para sua subsistência e que já vivem em condições de vulnerabilidade.

Na floresta amazônica, por exemplo, o aumento das queimadas, embora frequentemente associado a problemas socioeconômicos
e à negligência governamental, tem sido intensificado pelas mudanças climáticas. O aumento da temperatura e da área
suscetível ao fogo coloca em risco uma das regiões mais biodiversas do planeta e compromete a capacidade da Amazônia de
atuar como um importante sumidouro de carbono, essencial para mitigar as emissões globais. As populações nativas dessas
regiões, apesar de estarem entre as que menos contribuem para as emissões de carbono, são profundamente impactadas pelos
efeitos das queimadas, que aumentam a incidência de doenças respiratórias, e tem a sua mobilidade restringida por
mudanças no ciclo hidrológico que rege o transporte ribeirinho, o principal da região \cite{pinho}. A perda de
biodiversidade afeta diretamente a subsistência desses grupos, que dependem dos recursos naturais para sua alimentação e
sobrevivência. Como consequência da insegurança alimentar e da degradação de seus territórios, muitas comunidades acabam
migrando para centros populacionais próximos, onde enfrentam condições ainda mais vulneráveis, exacerbadas por novos
problemas socioeconômicos e pela dificuldade de adaptação ao ambiente urbano que frequentemente marginaliza suas
necessidades e tradições.

É a intrenseca vulnerabilidade de regiões como esta que torna mais importante as ações de prevenção e mitigação de
desastres naturais causados pelo aumento climático. Pois, mesmo que com a tecnologia atual não se possa reverter o
aumento de temperatura, é possível mitigar os efeitos derivados disto.

% Desenvolvimento sustentável; -> Vulnerabilidade dos que contribuem menos. Tentar ao máximo ficar no tópico de América
% Latina;

% IPCC AR6 WGI Chapter 11.4 (Heavy Precipitation) and Chapter 11.5 (Floods); TODO IMAGEM!
\subsection{Mudanças nos padrões de precipitação}
\label{rain}
Mudanças em média de temperatura são processos graduais que se tornam difícies de perceber em espaços curtos de tempo. A
comparação de temperatura entre um ano e o anterior não mostrará a tendência real das mudanças, especialmente
considerando efeitos naturais de variabilidade climática, como o El Niño e a La Niña. No entanto, essas mudanças
graduais afetam diretamente nos eventos climáticos extremos cuja visibilidade é mais clara. Esse aumento da temperatura,
coloca em desequilibrio delicados mecanismos atmosfericos levando à mudanças em, por exemplo, aumento de temperaturas
máximas, redução de temperaturas mínimas e ondas de calor\footnote{Definido por Dunn et al. (2024) \cite{dunn} como 3 ou
mais dias no qual a temperatura se encontra acima do 90º percentil}.

Um evento climático cujo mudança se destaca é a precipitação. Em essência, o aumento de temperatura leva à uma maior
evaporação da água e também a maior capacidade atmosférica de reter esse vapor de água, o que resulta em chuvas mais
intensas e irregulares. Uma definição mais completa pode ser dada por Seneviratne et al. (2021) no capítulo 11.1 do
sexto relatório, segundo grupo, do IPCC (AR6 WGII): 

\begin{citacao}
As mudanças na temperatura também controlam as variações no vapor d'água por meio do aumento da evaporação e da
capacidade de retenção de água da atmosfera. Em escala global, o conteúdo de vapor d'água integrado na coluna
atmosférica aumenta aproximadamente de acordo com a relação de Clausius–Clapeyron, com um aumento de cerca de 7\% para
cada 1°C de aquecimento médio global da superfície. (...) Projeções de modelos climáticos mostram que o aumento do vapor
d'água resulta em elevações significativas nos extremos de precipitação em todas as regiões, com uma magnitude que varia
entre 4\% e 8\% por 1°C de aquecimento da superfície.

\textit{Seneviratne et al. (2021)\cite{ipcc_ar6_wg2_chap11}, traduzido pelo autor.}

\end{citacao}

(EXPANDIR - Definir  melhor alta precipitação e o aumento dela);

É importante notar que o aumento das temperaturas não está somente associado a um incremento na precipitação total, mas
também à intensificação dos extremos de precipitação. Observa-se não só um aumento no volume máximo de chuva em um único
dia, evidenciando a maior intensidade dos eventos, como também no acumulado de precipitação em períodos de cinco dias,
indicando uma maior duração desses extremos \cite{dunn}.

Apesar da tendência global de aumento nas precipitações extremas, isso não significa que todas as regiões do planeta
experimentarão esse efeito de maneira uniforme. O aumento das chuvas vem acompanhado de mudanças nos padrões de
precipitação, influenciadas pelos mecanismos atmosféricos e pelas características geográficas locais. Assim, o
aquecimento global não apenas eleva o volume de chuvas, mas também altera sua distribuição, afetando onde e com que
frequência ocorrem.

\begin{figure}[htb] 
	\caption{\label{rain_tendency} Tendência de aumento de chuvas em diferentes regiões}
	\begin{center}
	    \includegraphics[scale=0.235]{img/SRL-image-0 (copy).png}
	\end{center}
	\legend{Fonte: \textit{IPCC Sixth Assessment Report - Working Group 1} \cite{ipcc_summary}}
\end{figure}

(EXPANDIR - Mencionar a imagem acima no texto explicando a região SES como alvo);

Essas variações são determinadas por particularidades climáticas regionais, o que dificulta a definição de uma tendência
uniforme. As mudanças nos padrões de precipitação estão diretamente relacionadas à um aumento na quantidade de secas e
de enchentes, pois, se a distribuição de precipitações aumenta ou diminui, os limiares que levam à esses eventos são
mais facilmente alcançados.

No Brasil, por exemplo, esses dois efeitos são observados de forma simultânea. No Nordeste, a região semiárida mais
populosa do mundo, há uma redução na frequência e intensidade das chuvas, o que tem levado a um aumento na ocorrência de
secas mais severas e frequentes, um problema crônico na área. Além disso, projeções indicam que o Nordeste enfrentará
aumentos de temperatura mais intensos do que outras regiões \cite{hoegh-guldberg}, agravando ainda mais os impactos da
seca. Por outro lado, na região Sudeste, observa-se um aumento nas chuvas intensas desde o início do século XX
\cite{dunn}, e, acompanhado das mesmas, uma intensificação em enchentes e deslizamentos, que são a maior causa de morte
relacionada à desastres naturais na região \cite{ufsc}. 
|
    \section{Enchentes}
    As enchentes estão entre os desastres naturais mais frequentes e impactantes do mundo, afetando milhões de pessoas e
causando prejuízos econômicos e sociais significativos todos os anos. Esses eventos, caracterizados pelo transbordamento
de corpos d’água ou pelo acúmulo excessivo de água em áreas urbanas e rurais, são resultado de uma complexa interação
entre fatores naturais e atividades humanas. No contexto das mudanças climáticas, a frequência e a intensidade das
enchentes têm aumentado em diversas regiões, tornando-se um desafio global que exige atenção urgente.

Neste capítulo, será abordado as principais características das enchentes, com foco especial nas enchentes fluviais, que
ocorrem quando rios transbordam. Como estudo de caso destes eventos, serão trazidos eventos recentes de enchentes
catastróficas e os seus impactos, um exemplo emblemático de como a combinação de mudanças climáticas, alterações na
morfologia dos rios e vulnerabilidade urbana pode levar a desastres de grandes proporções. Além disso, serão explorados
os processos hidrodinâmicos que influenciam a ocorrência e a duração das enchentes, bem como a intrínseca
vulnerabilidade das áreas urbanas a esses eventos extremos.

% ######################################################################################################################
% IPCC AR6 WGI - Chapter 8
\subsection{Enchentes e suas características}

\begin{figure}[htb] 
	\caption{\label{cicle} Ciclo da água}
	\begin{center}
	    \includegraphics[scale=0.165]{img/WaterCicle (copy).png}
	\end{center}
	\legend{Fonte: \textit{United States Geological Survey} \cite{usgs}, traduzido pelo autor}
\end{figure}

Enchentes são caracterizadas pelo aumento do nível da água em áreas que normalmente permanecem secas. Esse acúmulo
prolongado ocorre devido à dificuldade de escoamento, seja pela incapacidade do solo de absorvê-la ou pela lentidão na
drenagem natural, como o fluxo em grandes corpos d'água. Em geral, elas podem ser classificadas em três tipos
principais: enchentes pluviais, resultantes de chuvas intensas que excedem a capacidade de absorção do solo, causando
alagamentos; enchentes fluviais, provocadas pelo transbordamento de corpos d'água – geralmente rios – quando o nível da
água ultrapassa suas margens, frequentemente devido a chuvas persistentes, derretimento de gelo ou acúmulo de
sedimentos; e enchentes costeiras, que ocorrem em áreas litorâneas devido ao aumento do nível do mar, grandes ondas ou
fenômenos como ciclones e tsunamis.

As causas de uma enchente nem sempre podem ser facilmente atribuídas a um único evento, pois geralmente dependem de uma
combinação de fatores. No caso das enchentes fluviais, por exemplo, não é apenas a ocorrência de chuvas intensas no
corpo d'água principal ou em seus tributários que provoca o transbordamento. Se, por exemplo, o corpo hídrico tiver
capacidade de escoar a água de forma eficiente, o alagamento pode ser evitado. Muitos corpos d'água, cuja principal
fonte de abastecimento são enchentes pluviais, possuem uma morfologia adaptada para lidar com essa condição. Outros
fatores que contribuem para o transbordamento incluem a área de escoamento disponível em canais naturais e artificiais,
a umidade prévia do solo, sua capacidade de percolação, a presença de ventos fortes e sua interação com o fluxo da água,
a taxa de evaporação após o evento, a capacidade de absorção da vegetação local, o declive do terreno em relação à área
de escoamento, entre outros. Assim, não se pode associar exclusivamente o aumento das chuvas intensas (\autoref{rain})
ao crescimento destes eventos extremos \cite{ipcc_ar6_wg2_chap11}.

Mesmo assim, as enchentes são um fenômeno natural e frequentemente característico de certas regiões. Áreas com invernos
rigorosos, marcados pela acumulação de neve, costumam enfrentar enchentes durante o verão \cite{atashi}, quando o rápido
derretimento do gelo transporta grandes volumes de água de uma região para outra, sustentando ecossistemas que dependem
desse influxo sazonal. Regiões influenciadas por monções — ventos sazonais que trazem para o interior um acúmulo de
umidade e precipitação marítima — também são marcadas por cheias em seus corpos d'água, essenciais para a agricultura
local, como ocorre nos tributários dos rios Ganges e Amarelo. De maneira semelhante, as cheias anuais regulares do Rio
Nilo são fundamentais para a fertilidade da região. 

Entretanto, mesmo sendo uma parte integral do ciclo da água, a presença de assentamentos humanos geralmente corresponde
com a existência de grandes corpos d'água que podem ser, dadas circunstâncias, vulneráveis à enchentes. Essas imundações
podem invadir esses locais de forma inesperada e veloz, colocando vidas em risco, comprometendo a segurança das
comunidades e causando danos estruturais que amplificam a disrupção das atividades humanas. Fatalidades causadas por
afogamentos são comuns em regiões alagadas, não apenas pela rapidez com que as enchentes ocorrem, mas também devido às
características do escoamento das águas, que podem gerar correntes fortes e imprevisíveis. As enchentes também criam
condições propícias para o agravamento de outras condições sanitárias. Elas podem transmitir doenças por meio da água
contaminada, além de comprometer sistemas críticos, como a contaminação de fontes de água potável, a interrupção do
fornecimento de energia elétrica, afetando hospitais e serviços essenciais, e a destruição de cadeias de suprimentos.
Esses fatores podem isolar comunidades, mesmo aquelas não diretamente afetadas pela inundação, dificultando o acesso a
mantimentos e recursos básicos.

As enchentes urbanas representam um dos principais vetores de risco em áreas densamente povoadas, combinando
vulnerabilidade socioambiental com impactos econômicos e humanos catastróficos. 

% Em Roma, por exemplo, exitem registros de cheias no Rio Tibre que remontam ao século 6 BCE, quando a construção da
% Cloaca Máxima marcou um dos primeiros sistemas organizados de gestão hidráulica urbana. Já em 1870, uma falha
% estrutural no canal do rio resultou em uma inundação de 17m acima do nível normal, submergindo monumentos icônicos
% como o Panteão, localizado nos Campos de Marte, uma região considerada como várzea.

% % Photo: Miss I. B. Trevenen, Rome. The Floods in Rome, the Pantheon during the Inundation. Illustration for The %
% Illustrated London News, 15 December 1900. English Photographer (20th Century)

% \begin{figure}[htb] \caption{\label{mars} Enchente de 1900 na entrada do Panteão} \begin{center}
%   \includegraphics[scale=0.55]{img/pantheon.jpg} \end{center} \legend{Fonte: \textit{The Illustrated London News}
%   \cite{london}} \end{figure}

Mais recentemente, destaca-se a enchente de 2003 na provincia de Santa Fe, na Argentina. Depois de dias de chuvas
intensas, o rio Salado transbordou no dia 28 de Abril. Aproximadamente 100 mil pessoas tiveram de ser evacuadas da
região e 160 faleceram. As perdas estimadas da região ficaram na faixa de 500 milhões de pesos argentinos
\cite{santa_fe}.

% Os efeitos das enchentes, no entanto, não se limitam às comunidades urbanas. Atividades econômicas fundamentais, como
% a agricultura, estão entre as mais prejudicadas por esses desastres. A necessidade de grandes extensões de terra e a
% proximidade de corpos d’água para irrigação tornam o setor agrícola particularmente vulnerável. Em situações de
% enchente, a safra pode ser destruída não somente pela força da água, mas também pela permanência prolongada de água
% nessas regiões, já que muitas espécies de plantas não conseguem sobreviver a tais condições. Além da destruição de
% colheitas e da perda de gado, as enchentes frequentemente resultam em danos à infraestrutura essencial, como sistemas
% de irrigação, equipamentos, instalações de processamento e transporte, além de alterações nas terras cultiváveis. De
% acordo com a Organização de Alimentos e Agricultura das Nações Unidas \cite{onu}, enchentes são responsáveis por mais
% da metade de todas as perdas em colheitas causadas por desastres naturais. Entre 2008 e 2018, essas perdas em países
% em desenvolvimento ultrapassaram 21 bilhões de dólares, tornando as enchentes a segunda maior causa de danos
% econômicos no setor agrícola, atrás apenas das secas. Em países em desenvolvimento, onde setores primários têm um peso
% significativo na economia, os efeitos desses desastres são amplamente sentidos de forma indireta no enfraquecimento
% econômico, podendo levar a uma maior situação de vulnerabilidade e despreparo.

% O setor agrícola de Bangladesh, onde 60\% da população depende diretamente da agricultura, é recurrentemente afetado
% por enchentes. Em 2020, inundações no nordeste do país destruíram 220.000 hectares de plantações de arroz, afetando
% 4,7 milhões de pessoas. A FAO estima que, entre 2015 e 2022, perdas médias anuais chegaram a US\$ 1,2 bilhão,
% equivalentes a 2,3\% do PIB agrícola nacional. Esses eventos não apenas comprometem a segurança alimentar, mas também
% perpetuam ciclos de pobreza, já que pequenos produtores perdem até 80\% de sua renda sazonal.

% ######################################################################################################################
\subsection{Tendências de enchentes devido à mudanças climáticas}
De maneira direta e indireta, as atividades humanas alteram profundamente o ciclo hidrológico. Conforme demonstrado na
\autoref{rain}, o aquecimento global não apenas eleva as temperaturas, mas também intensifica a frequência e a magnitude
das chuvas, impulsionado pela maior quantidade de umidade presente na atmosfera. Embora as variações naturais possam
dificultar a identificação de mudanças locais nos eventos de chuva, há evidências de que, quando padrões meteorológicos
naturais promovem chuvas prolongadas em um clima aquecido, a intensidade desses episódios é amplificada pela abundância
de vapor d’água. Isso resulta em um aumento de eventos extremos, mesmo que regionalmente as mudanças possam se
manifestar de formas diversas – com áreas áridas tornando-se ainda mais secas e regiões úmidas sendo afetadas por chuvas
mais intensas.

Entretanto, é importante destacar que chuvas mais intensas nem sempre se traduzem automaticamente em enchentes mais
graves, pois os impactos dependem de múltiplos fatores, como as características da bacia hidrográfica, a topografia, a
duração do evento e o estado de umidade do solo antes da ocorrência da chuva. Além disso, a gestão do uso do solo – como
o desmatamento para expansão agrícola ou a urbanização desordenada – pode acelerar o fluxo de água, direcionando-a de
forma mais rápida para áreas vulneráveis. Em contrapartida, a extração intensiva de água dos rios pode, em certos
contextos, reduzir os níveis dos cursos d’água e, temporariamente, mitigar o risco de enchentes.

Portanto, mesmo que as variações naturais dificultem atribuir as mudanças locais unicamente ao aquecimento global,
existe uma tendência de que o aumento da concentração de gases de efeito estufa contribuem para chuvas diárias cada vez
mais intensas e eventos extremos. Esse cenário, aliado às alterações no manejo do solo e das bacias hidrográficas, leva
à enchentes mais severas em muitas regiões.

Na região Sudeste do Brasil, as tendências de enchentes fluviais se mostram de acordo com os padrões globais. Como é uma
área caracterizada por longos períodos úmidos, espera-se que esses períodos se tornem ainda mais frequentes e intensos.
Em regiões com alta modificação antropogênica, especialmente devido à atividade agrícola, esse agravamento tende a ser
ainda mais pronunciado \cite{chagas}. Eventos de enchentes catastróficas, tradicionalmente calculados para ocorrer uma
vez por século, já apresentam indícios de um período de retorno menor, levando à eventos recordes mais frequentes. De
acordo com a análise de extremos hidrológicos de Vieira et al. (2022) \cite{vieira}, a probabilidade de ocorrência de
enchentes devido à precipitação e fluxo extremo na Bacia Hidrográfica do Guaíba é maior e pode ultrapaçar marcos
históricos, como ocorreu com as enchentes de 2024.

% ######################################################################################################################
\subsection{Morfologia de corpos fluviais e Hidrodinâmica de enchentes}
\label{morf}
Um rio é uma formação hidrológica que transporta água e sedimentos de regiões mais elevadas, conhecidas como cabeceiras,
para áreas de menor altitude em direção à corpos maiores de água, seguindo um caminho definido, denominado canal. Esses
canais desempenham um papel fundamental no ciclo hidrológico, conduzindo a maior parte da água das regiões interioranas
até as áreas costeiras. Ao longo de seu trajeto, os rios erodem o solo e as rochas, transportando sedimentos que,
eventualmente, são depositados em lagos e oceanos \cite{holden}.

% \begin{figure}[htb] \caption{\label{guaiba} Rio Guaíba} \begin{center}
%   \includegraphics[scale=0.113]{img/IMG_20250406_163317.jpg} \end{center} \legend{Fonte: Reprodução própria}
%   \end{figure}

A formação dos canais fluviais depende principalmente de fatores como precipitação, derretimento de neve e água
subterrânea. A precipitação, em particular, destaca-se como um fator crucial. Além de contribuir diretamente para o
abastecimento dos rios, ela promove a infiltração de água no solo, originando aquíferos e água subterrânea. Grande parte
dessa água subterrânea acaba integrando a bacia hidrográfica, também conhecido como sistema de drenagem, por meio de
riachos e córregos, que se combinam para formar canais maiores. Esse processo inicia-se com o escoamento superficial,
onde a água da chuva flui inicialmente como uma fina camada sobre o solo, que se concentra em áreas de maior declive e
começa um processo de erosão do solo, formando pequenos sulcos. Com o tempo, esses sulcos se aprofundam e se fundem,
dando origem a canais mais amplos que estabelecem uma rede de drenagem, formando afluentes que se unem a um rio
principal, criando-se uma complexa rede de drenagem \cite{nelson}.

As bacias hidrográficas podem assumir diversas formas, dependendo das características geográficas e geológicas da região
onde a água escoa. O principal fator é o tipo de rocha presente na área, pois diferentes rochas apresentam resistências
variadas à erosão. Além disso, a topografia local, como a inclinação do terreno e a presença de falhas geológicas,
também ajudam a definir esses sistemas. Dessa forma, o padrão de drenagem de uma região pode ser um indicativo de sua
geologia e topografia subjacentes. Na \autoref{drain}, são apresentados alguns dos tipos mais comuns de bacias
hidrográficas, cada um associado a características específicas do terreno.

\begin{figure}[htb]
	\caption{\label{drain} Exemplos de formas de bacias hidrográficas}
	\begin{center}
	    \includegraphics[scale=0.60]{img/RiverShapes.png}
	\end{center}
	\legend{Fonte: Reprodução própria}
\end{figure}

Na \autoref{drain}.1, observa-se o sistema de drenagem dendrítico, cujo padrão se assemelha aos ramos de uma árvore. É
típico de áreas compostas por granitos ou arenitos, onde a resistência à erosão é uniforme, permitindo o escoamento em
todas as direções, geralmente em planícies ou relevos suaves. A \autoref{drain}.2 mostra o sistema radial, em que os
cursos d’água divergem de um ponto elevado central, comum em regiões montanhosas com forte variação altimétrica. Já a
\autoref{drain}.3 exibe o padrão retangular, com canais angulados, comuns em áreas com falhas ou fraturas na rocha, onde
os rios seguem zonas com diferentes resistências à erosão.

% Esses padrões de drenagem influenciam diretamente a morfologia e a dinâmica dos rios ao longo de seu curso. Nas áreas
% iniciais do canal, próximas à nascente, o gradiente é mais acentuado, refletindo a maior elevação e a forte tendência
% da água de escoar em uma direção específica. À medida que o rio avança em direção à foz, o gradiente diminui, e o
% fluxo se torna mais lento. Nessas regiões de baixo gradiente, é comum a deposição de sedimentos que leva à formação de
% várzeas, ecossistemas úmidos alimentados pelas inundações periódicas do rio, e meandros, curvas sinuosas que resultam
% da erosão nas margens côncavas e da deposição de sedimentos nas margens convexas. Meandros mais acentuados são
% característicos de rios mais velhos e bem estabelecidos. As várzeas, por sua vez, podem servir como uma barreira para
% as enchentes, dado que a sua disposição plana situada adjacentemente aos rios favorece a absorção do excesso de água
% durante cheias, mas também são vulneráveis a inundações prolongadas quando a capacidade de infiltração do solo é
% ultrapassada.

% Essa evolução longitudinal do rio também se manifesta na sua geometria, especialmente na seção transversal do canal. A
% seção transversal de um rio é definida como a área por onde ocorre o fluxo de água no canal. Essa área varia ao longo
% do curso do rio e reflete as características do fluxo, como velocidade, descarga e capacidade de transporte de
% sedimentos. As partes mais profundas do canal geralmente estão associadas às regiões onde a velocidade da corrente é
% maior, pois a água concentra sua energia em áreas específicas, promovendo maior erosão. Além disso, tanto a largura
% quanto a profundidade da seção transversal aumentam conforme a descarga do rio, já que um maior volume de água exige
% uma área mais ampla para escoar. Por fim, a forma da seção transversal costuma ser irregular e variável ao longo do
% trajeto, influenciada por fatores como a composição das rochas, a dinâmica de erosão e deposição de sedimentos, e a
% presença de obstáculos naturais ou artificiais. A segunda parte da \autoref{sinus} ilustra a seção transversal de um
% rio, destacando em especial a relação entre a forma do canal e a velocidade da água.

Os padrões de drenagem condicionam a evolução longitudinal dos rios. Nas cabeceiras, o elevado gradiente promove fluxos
rápidos e erosão. À medida que o gradiente diminui em direção à foz, a redução da velocidade favorece a deposição de
sedimentos, formando várzeas - planícies de inundação adjacentes que atuam como áreas de amortecimento durante cheias -
e meandros por erosão seletiva em margens côncavas e deposição nas convexas. Esta morfologia reflete-se igualmente na
seção transversal do canal, onde a geometria (largura, profundidade e rugosidade) determina a relação entre velocidade,
descarga e capacidade de transporte sedimentar. Como ilustrado na \autoref{sinus}, áreas de maior profundidade
correlacionam-se com velocidades elevadas e erosão, enquanto aumentos na descarga ampliam as dimensões do canal. A forma
transversal é dinamicamente modulada por fatores litológicos, processos erosivo-deposicionais e intervenções antrópicas.

\begin{figure}[htb]
	\caption{\label{floodplain} Exemplo de regiões de várzea inundadas no Rio Severn, Inglaterra}
	\begin{center}
	    \includegraphics[scale=0.35]{img/floodplain.jpg}
	\end{center}
	\legend{Fonte: Tom Blockley \cite{varzea}}
\end{figure}

A seção transversal do canal regula a relação entre largura, profundidade e velocidade do fluxo. Canais estreitos e
profundos concentram a água em alta velocidade, aumentando a capacidade de transporte de sedimentos, mas também elevando
o risco de transbordamento rápido caso a descarga exceda a capacidade de contenção. Em contraste, canais largos e rasos
dissipam parte da energia do fluxo, reduzindo a velocidade, porém exigindo maior área para acomodar volumes elevados de
água. Alterações na seção transversal, como estreitamentos naturais ou acumulo de sedimentos, podem amplificar a
profundidade do fluxo e a pressão sobre as margens, favorecendo transbordamentos.

Além disso, a velocidade da água influencia não apenas a erosão e a deposição, mas também a capacidade do rio de
transportar materiais em suspensão. Irregularidades no leito e nas margens aumentam a resistência ao fluxo, reduzindo a
velocidade da água próxima às margens e intensificando a erosão nessas áreas. Em contraste, a água no centro do canal é
menos afetada por esse efeito, mantendo uma velocidade mais elevada. Além disso, os efeitos da fricção são mais
pronunciados em canais largos e rasos, enquanto em canais estreitos e profundos a resistência é menor. Outro fator
crucial é o gradiente do canal: regiões com maior declive tendem a apresentar águas mais velozes, enquanto áreas de
baixo gradiente favorecem fluxos mais lentos.

A sinuosidade do canal também desempenha um papel importante na distribuição da velocidade da água. Em canais retos, a
velocidade máxima ocorre no centro do canal, devido à menor influência da fricção com as margens. Já em canais sinuosos,
a velocidade é mais alta na parte externa das curvas, onde a força centrífuga empurra a água contra as margens,
resultando em maior profundidade e erosão nesses pontos. Por outro lado, na parte interna das curvas, a velocidade é
reduzida, favorecendo a deposição de sedimentos transportados pela água. Rios excessivamente sinuosos tendem a retardar
o escoamento, prolongando a permanência de água na bacia durante chuvas intensas e aumentando a probabilidade de
inundações. Por outro lado, a retificação artificial de meandros, embora acelere o fluxo, pode transferir riscos de
enchentes para áreas antes menos expostas.

\begin{figure}[htb]
	\caption{\label{sinus} Relação entre sinuosidade e velocidade em diferentes pontos de um rio}
	\begin{center}
	    \includegraphics[scale=0.40]{img/RiverSinus.png}
	\end{center}
	\legend{Fonte: Reprodução própria}
\end{figure} 

A carga de sedimentos é outra característica importante na morfologia dos rios. Ela reflete o volume de material
transportado, variando conforme a energia do fluxo e as propriedades do leito. Esta carga engloba tanto minerais
dissolvidos, derivados da erosão rochosa e responsáveis pela salinização de corpos receptores, quanto partículas sólidas
em suspensão, como fragmentos de rochas e matéria orgânica. O tamanho e o peso dessas partículas transportadas dependem
diretamente da velocidade do rio, dado que em fluxos mais lentos, as partículas tendem a se depositar no leito,
reduzindo sua mobilidade.

A descarga, por sua vez, representa o volume de água que passa por uma seção do rio em um determinado tempo e está
ligada à larguera e à profundidade do canal. A mesma pode ser dada em metros cúbicos de água por segundo ($m^3/s$) e é
definida pela da seção transversal do rio ($m^2$) e a velocidade da água ($m/s$) através de $D = A*v$, onde $D$ é a
descarga, $A$, a seção transversal e $v$ a velocidade. A descarga de um rio pode aumentar devido a fatores como a
precipitação no corpo d'água principal e em seus afluentes, pela contribuição de água subterrânea que se integra ao
sistema da bacia hidrográfica, ou pela quebra de barreiras naturais e/ou artificias, como barragens formadas por
sedimentos e matéria orgânica. Esses aumentos repentinos da descarga são os principais indicadores de enchentes, pois
podem superar a capacidade de armazenamento temporário do canal e das várzeas.

A relação entre descarga e enchentes é direta: sempre que a vazão excede a capacidade do canal, as águas transbordam
para as várzeas adjacentes e, em eventos extremos, para áreas além dessas planícies. Esses picos estão diretamente
associados à intensidade e duração das chuvas (\autoref{rain}), mas também ao estado prévio de umidade da bacia
hidrográfica. Em bacias já saturadas por chuvas anteriores, a água da precipitação escoa diretamente para os canais,
amplificando a descarga. Em contraste, solos secos absorvem grande parte da água, retardando o impacto no fluxo fluvial.

A capacidade de absorção do solo, determinada pela cobertura vegetal e pela porosidade – característica que determina a
taxa de infiltração da água e sua percolação\footnote{Movimento da água através das camadas do solo, favorecendo o
armazenamento em profundidade}, atua como um regulador natural. Solos porosos e ricos em matéria orgânica permitem maior
infiltração e percolação da água, enquanto compactação, urbanização ou camadas impermeáveis reduzem essa capacidade.
Quando o limite de saturação é atingido, o escoamento superficial direciona-se para canais menores, cuja sincronização
em múltiplos afluentes gera uma descarga acumulada crítica no rio principal, amplificando o risco de inundações. Logo,
um solo saturado e com pouca capacidade de inflitração resulta em enchentes mais frequentes em comparação com um solo
seco que absorve a água das chuvas sem exceder a capacidade dos canais.

A hidrodinâmica de enchentes fluviais resulta da interação complexa entre múltiplos fatores geomorfológicos e
hidrológicos. O transbordamento de corpos hídricos nunca decorre de um único elemento isolado, mas sim da interação
entre processos que, em condições críticas, potencializam eventos extremos. Podemos sintetizar esse mecanismo em três
estágios consecutivos: a ocorrência de precipitações intensas e prolongadas que excedem a capacidade de infiltração do
solo; o consequente escoamento superficial que eleva abruptamente a descarga hídrica; e a superação da capacidade de
vazão definida pela seção transversal do canal, levando à inundação das áreas adjacentes - particularmente as várzeas.
Esse processo evidencia como características físicas do canal, condições antecedentes do solo e padrões climáticos atuam
de forma integrada na gênese das enchentes.

% ######################################################################################################################
\subsection{Enchentes em ambientes urbanos e sistemas de prevenção}
Centros urbanos, conglomerados socioeconômicos críticos consolidados com frequente déficit de planejamento, operam
através de redes assimétricas de serviços essenciais (água, energia, transporte) cuja interdependência gera
vulnerabilidade sistêmica, com onde falhas pontuais levando à efeitos em cascata. Este quadro agrava-se diante das
projeções demográficas, pois, enquanto é estimado que a população humana aumente de 7.94 bilhões para 9.68 bilhões entre
o período de 2024 até 2050, (21\% em 26 anos), também é estimado que a população em centros urbanos cresça de 4.61
bilhões para 6.7 billhões durante 2024-2050, i.e., um crescimento de 30\% no mesmo período
\cite{World_Population_Prospects}. Tal expansão acelerada converterá áreas verdes (notadamente várzeas) em superfícies
impermeáveis, reduzindo a capacidade de infiltração do solo e ampliando o escoamento superficial pluvial.
Consequentemente, milhões de novos habitantes concentrar-se-ão em zonas com resiliência hidrológica comprometida,
expondo-os a riscos de enchentes superiores aos períodos pré-urbanização \cite{hammond}.

% \begin{figure}[htb] \caption{\label{katmandu} Enchente urbana em Kathmandu} \begin{center}
%   \includegraphics[scale=0.65]{img/kapan-flood-chandra.jpg} \end{center} \legend{Fonte: Chandra Bahadur Aale
%   \cite{chandra}} \end{figure}

As cidades enfrentam elevado risco de enchentes fluviais exatamente pelas alterações promovidas nessas zonas – comuns em
grandes centros urbanos devido à proximidade histórica de rios, dada a necessidade de acesso a água potável e descarte
de efluentes. Logo, as várzeas, outrora absorventes, tornam-se epicentros de risco. Quando transbordam, os rios não
apenas inundam essas planícies urbanizadas, mas também comprometem sistemas interdependentes, como sistemas de
distribuição de água, esgotos e infraestruturas de energia que sofrem danos em cascata, gerando impactos em regiões
distantes do epicentro do transbordamento.

A capacidade de eventos extremos causarem impactos significativos pode ser descrita através de três fatores: risco,
exposição e vulnerabilidade. Para que centros urbanos possam construir resiliência, é necessário primeiro entender estes
fatores. Desta forma, o planejamento urbano passa a ser menos reativo à situações de crise e mais proativo na sua
construção de adaptabilidade, levando à menores perdas de vidas em curto prazo e recuperação mais eficiente à longo
prazo, visando a preservação e restauração da funcionalidade de estruturas essênciais que compõem as malhas urbanas em
um ambiente em mudança. 

O risco, definido pela probabilidade de ocorrência do evento, está intrinsecamente vinculado aos fatores geográficos,
como padrões pluviométricos, e meteorológicos, como características topográficas, inerentes à localização do centro
urbano, assim como foi descrito na \autoref{morf}. Por depender destas variáveis naturais, trata-se do elemento sobre o
qual há menor controle direto por parte dos gestores urbanos, pois embora medidas governamentais possam ser tomadas
vizando a redução da crise climática, as mesmas são ineficazes quando aplicadas somente em escala regional.

% \begin{figure}[htb] \caption{\label{risco} Relação entre risco, exposição e vulnerabilidade em sistemas urbanos}
%   \begin{center} \includegraphics[scale=0.274]{img/RiverRisk.png} \end{center} \legend{Fonte: Bruno Merz et. al (2021)
%   \cite{merzcauses}} \end{figure}

A exposição, por outro lado, refere-se aos elementos materiais e econômicos suscetíveis a danos durante uma enchente,
sendo parcialmente controláveis por meio de políticas de planejamento urbano, como o zoneamento. Essa dimensão é
frequentemente quantificada pelo nível da água, que permite estimar quais ativos (infraestruturas, edificações, redes de
serviços) serão afetados em diferentes cenários de inundação. Estudos como o de Merz et al. \cite{merz} destacam,
entretanto, que essa métrica não pode explicar a totalidade dos danos, dado que fatores hidrodinâmicos adicionais, como
velocidade da corrente e carga de sedimentos, desempenham papéis críticos.

Apesar dessas nuances, o nível da água mantém utilidade como indicador generalizado de impacto, dado que os outros
fatores são altamente dependentes do tipo de estrutura afetado, e, logo, o mesmo oferece uma base objetiva para
estratégias de zoneamento urbano que visem minimizar perdas. Além disso, o monitoramento do nível da água representa uma
forma unificada de entender dinâmicas complexas de inundação. Embora o conhecimento do mesmo só seja efetivo caso
aplicado em mapas de inundação, a redução da complexidade para um único valor facilita a comunicação do risco para a
população e agiliza a tomada de decisão baseado nas capacidade das estruturas de defesa presentes. Neste contexto,
sistemas de monitoramento contínuo do nível d'água assumem papel central, pois fornecem dados em tempo real que
alimentam modelos preditivos e sistemas de alerta precoce.

A vulnerabilidade, por sua vez, é o fator de maior controle para tomadores de decisões de centros urbanos. Estas são as
medidas de adaptabilidade tanto de estruturas, quanto de habitações perante eventos extremos. As medidas relacionadas à
redução da vulnerabilidade incluem desde infraestruturas físicas, como sistemas de defesa contra as enchentes, sistemas
de prevenção, até sistemas de gerenciamento de crises e alerta de riscos. 

Em escala regional, barragens destacam-se como sistemas estratégicos de armazenamento hídrico, capazes de regular fluxos
extremos provenientes de chuvas intensas. Além da geração de energia, as represas podem alocar partes dos seus
reservatórios para contenção de cheias, liberando água de forma gradual no curso principal da bacia hidrográfica de
acordo com a capacidade da mesma. Essa modulação reduz a velocidade de escoamento, minimizando danos por erosão e
transbordamentos.

% \begin{figure}[htb] \caption{\label{itaipu} Usina hidroelétrica de Itaipu com seus vertedouros abertos} \begin{center}
%   \includegraphics[scale=0.12]{img/Itaipu_geral.jpg} \end{center} \legend{Fonte: Jonas de Carvalho \cite{itaipu}}
%   \end{figure}

Uma solução mais comum em centros urbanos se dá nas formas do redirecionamento dos cursos naturais dos rios, em especial
com a retificação de canais naturais substituindo a sinuosidade fluvial por trajetos retilíneos artificialmente
construídos, dado que canais sinuosos diminuem a velocidade as águas em seus meandros e retardam o escoamento (vide
\autoref{sinus}). Os canais retificados não apenas direcionam a vazão crítica para corpos receptores com menor risco de
transbordamento, mas também integram-se à infraestrutura urbana como componentes de sistemas de abastecimento hídrico e
monitoramento de condições fluviais. Além disso, eles podem contar com sistemas de desvio que visam ligar partes
distintas da rede hidrográfica, desviando parcelas da vazão para rotas alternativas que contornam núcleos urbanos
densamente povoados.

Como alternativa complementar à retificação de canais, o aterramento fluvial amplia a capacidade de contenção hídrica
mediante estruturas de expansão das margens – como diques ou aterros. Essas intervenções aumentam a seção transversal do
canal, aumentando sua capacidade de vazão antes do transbordamento. Contudo, a eficácia de ambas estas estruturas
depende criticamente da integridade estrutural: rupturas pontuais, como brechas em diques ou falhas de vedação,
convertem essas barreiras em focos de amplificação hidráulica. Nessas situações, a súbita liberação de energia cinética
concentrada em áreas adjacentes, muitas vezes desprovidas de drenagem redundante e com falta de aviso prévio, exacerba
inundações locais.

Além das intervenções superficiais, sistemas subterrâneos de drenagem pluvial complementam a gestão hidráulica urbana.
Projetadas para transportar água separadamente do esgoto, as galerias pluviais operam como redes de canais subterrâneos
que direciona o escoamento superficial para canais, sejam eles artificiais ou naturais, ou para estruturas de contenção.

% -> Essas galerias pluviais podem até mesmo apresentar riscos maiores quando levado em consideração a sua habitação por
% moradores de rua, como nos túneis de Las Vegas. Em situações de chuvas intensas, nas quais os túneis serviriam o seu
% propósito de escoamento alternativo, essa parcela da população se encontra em um enorme risco de afogamento, visto que
% o fluxo da água nessas estruturas é súbito e intenso.

Integradas a essas redes, estações de bombeamento hidráulico atuam na transferência controlada de água entre bacias,
equalizando cargas em áreas de gradiente topográfico desfavorável. Em escala subterrânea, reservatórios de amortecimento
operam como análogos submersos de lagos artificiais, armazenando volumes críticos em períodos de pico para posterior
liberação gradual via sistemas de drenagem. Esses reservatórios permitem tanto a descarga regulada para canais
artificiais quanto a inflitração direta do lençol freático, contornando a impermeabilidade do solo urbano através destes
processos de percolação assistida. Essa dualidade funcional transforma a infraestrutura em mecanismo híbrido de controle
e restauração hidrológica.

Os reservatórios de retenção seca representam soluções multifuncional que operam como espaços públicos ordinários e
bacias de acumulação durante eventos extremos. Localizados estrategicamente em áreas de baixa cota topográfica (como
parques e praças), aproveitam o gradiente natural para interceptar fluxos superficiais, reduzindo a carga sobre redes de
drenagem em períodos chuvosos enquanto mantêm funções recreativas em condições normais. Sua integração urbanística
minimiza impactos visuais e sociais, com zonas de infraestrutura de baixa criticalidade que garantem rápida recuperação
operacional pós-inundações - característica particularmente valiosa em áreas densamente ocupadas com espaço limitado. O
Japão, devido à sua alta densidade urbana, exemplifica esta abordagem no complexo Parque Shin-Yokohama/Estádio Nissan
\cite{shinyokohamaplan}, projetado para conter enchentes do rio Tsurumi adjacente (\autoref{shin_yoko}). Durante cheias,
a água é direcionada ao parque por diques (\autoref{shin_yoko}.1), armazenada temporariamente (\autoref{shin_yoko}.2) e
liberada controladamente quando o canal recupera capacidade (\autoref{shin_yoko}.3), complementado por alterações no
canal fluvial e elevação estrutural do estádio para garantir funcionalidade durante eventos extremos.

\begin{figure}[htb]
	\caption{\label{shin_yoko} Parque de Shin-Yokohama}
	\begin{center}
	    \includegraphics[scale=2.3]{img/Lago-parque-japao.png}
	\end{center}
	\legend{Fonte: Governo da Prefeitura de Kanagawa \cite{shinyokohamapark} e Ministério de Terras, Infraestrutura,
	Transporte e Turismo, Escritório de Desenvolvimento Regional de Kanto \cite{shinyokohamaplan}}
\end{figure}


% Shin-Yokohama Park (新横浜公園): https://www.nissan-stadium.jp/csr/safety03.php Estádio Nissan (日産スタジアム):
% https://nissan-stadium.j-server.com/LUCNISSANS/ns/tl.cgi/https://www.nissan-stadium.jp/csr/safety03.php
% https://www.ktr.mlit.go.jp/keihin/keihin00119.html
% https://www-ktr-mlit-go-jp.translate.goog/keihin/keihin00119.html?_x_tr_sl=auto&_x_tr_tl=en&_x_tr_hl=en-US

Enquanto as Infraestruturas Cinza, grandes sistemas de engenharia convencional em concreto para controle direto de
fluxos, representam soluções predominantes, as Infraestruturas Verdes emergem como abordagem complementar baseada em
processos naturais. Estas utilizam componentes vegetados e permeáveis (parques, jardins, telhados vivos e rearborização
de vias) para substituir superfícies impermeáveis, aumentando a capacidade de absorção urbana. Sua implementação em
zonas industriais, comerciais e residenciais visa não apenas a gestão sustentável de águas pluviais, mas também a
mitigação de desequilíbrios urbanos, proporcionando melhorias da qualidade do ar, redução de ilhas de calor e incremento
da biodiversidade. Esta multifuncionalidade torna-as particularmente valiosas em áreas de alta densidade, contrastando
com as limitações das estruturas cinza: elevados custos de manutenção, potencial deslocamento de riscos para áreas
adjacentes e menor resiliência ecológica. Quando integradas como sistemas complementares, as soluções verdes oferecem
suporte sustentável à redução de riscos de enchentes urbanas \cite{gill}.

% A mitigação de riscos hidrológicos urbanos exige a complementaridade entre medidas estruturais e não estruturais,
% sendo estas últimas fundamentais para reduzir a vulnerabilidade sistêmica. Medidas de zoneamento emergem como
% estratégia central, delimitando restrições à ocupação em planícies aluviais mediante modelagem hidrodinâmica de
% cenários de enchente. A relocação preventiva de infraestruturas críticas nessas zonas – embora politicamente complexa
% – reduz exponencialmente a exposição de ativos estratégicos, enquanto a reconversão de áreas ocupadas em zonas verdes,
% como parques ou áreas públicas. Além disso, para infraestruturas que não podem ser removidas destas áreas, é essencial
% a adaptação das mesmas buscando mais resiliência.

A eficácia dessas ações depende criticamente de sistemas preditivos integrados, combinando redes de monitoramento
hidrometeorológico, como estações fluviométricas e pluviométricas para monitoramento de níveis de rios e eventos
extremos, com plataformas de análise em tempo real. Tais sistemas permitem a geração de alertas precoces baseados em
limiares de precipitação acumulada e níveis hidrométricos críticos. Contudo, o ciclo de gestão de risco só se completa
com protocolos de comunicação efetiva, que convertam estes dados técnicos em informações acionáveis para a população.
Sistemas de alerta georreferenciados, como a implementação de sirenes moduladas, avisos via mensagens de SMS e APIs para
integração em redes sociais, devem ser continuamente adaptáveis para garantir a comunicação efetiva dos riscos e
aumentar a segurança comunitária.

Um exemplo de sistema de monitoramento integrado nacional vem através do Centro Nacional de Monitoramento e Alertas de
Desastres Naturais (CEMADEN) que opera uma rede de monitoramento hidrometeorológico no Brasil. Seu modelo combina coleta
de dados in situ através de sensores pluviométricos e fluviométricos instalados em áreas de risco, com análises e
imagens de satélite \cite{cemadem}. Essa infraestrutura permite a emissão de alertas antecipados para eventos extremos
(como enchentes e deslizamentos), direcionados a defesas civis municipais e estaduais. Contudo, desafios operacionais
persistem, incluindo falhas na atualização de limiares críticos e lacunas na densidade de cobertura instrumental, como
evidenciado durante as enchentes de 2024 no Rio Grande do Sul, onde limitações na transmissão de alertas impactaram a
eficácia das respostas emergenciais.
    \section{As Enchentes no Rio Grande do Sul}
    \subsection{Morfologia do Rio Guaíba}
% Secretaria do Meio Ambiente e Infraestrutura: https://sema.rs.gov.br/g080-bh-guaiba

Localizada na região metropolitana do Rio Grande do Sul, a Bacia Hidrográfica do Guaíba constitui uma das Bacias mais
importantes do estado, possuíndo uma população de 1.3 milhões de pessoas nas cidades que percorrem o seu curso. Entre
elas se destacam os municípios da Região metropolitana de Porto Alegre (RMPOA), como Porto Alegre, Eldorado do Sul,
Guaíba e Canoas. 

\begin{figure}[htb]
	\caption{\label{guaiba_pol} Mapa Político do Guaíba}
	\begin{center}
	    \includegraphics[scale=0.48]{img/guaiba_pol.png}
	\end{center}
	\legend{Fonte: Instituto Brasileiro de Geografia e Estatística \cite{ibgeguai}}
\end{figure}


A mesma é principalmente formada pelas águas dos rios que desembocam no Guaíba, em especial aqueles que formam o Delta
do Jacuí, incluindo o Rio Caí, o Rio dos Sinos, o Rio Gravataí e o próprio Rio Jacuí. A região do Delta em si conta com
centros urbanos das cidades de Guaíba e Eldorado do Sul, além de contar com diversas ilhas, como a Ilha das Flores, a
Ilha Grande dos Marinheiros e a Ilha da Pintada. Alguns bairros de Porto Alegre também parte do sistema do ilhas, como o
Humaitá e o Navegantes.

% \begin{figure}[htb] \caption{\label{bacia_guaiba} Mapa da Bacia Hidrográfica do Lago Guaíba} \begin{center}
%   \includegraphics[scale=0.66]{img/bacia_guaiba.png} \end{center} \legend{Fonte: Secretaria do Meio Ambiente e
%   Infraestrutura do Rio Grande do Sul \cite{guaiba_bacia}} \end{figure}

Situa-se como destino das águas do Delta do Jacuí, destas, a descarga é composta primariamente pelo Rio Jacuí (84,6\%),
com o Rio dos Sinos, 7,5\%, Rio do Caí, com 5,2\%, e Rio Gravataí, com 2,7\% formando o restante da vazão. O Guaíba em
si pode ser tanto classificado como um Rio, quanto um lago. Historicamente, o mesmo já foi classificado de ambas as
formas, além de ria e estuário. O mesmo representa um ambiente de transição, com características de ambos rio e lago,
como um curso fluvial e uma dinâmica de armazenamento lacustre. Neste estudo, adota-se a nomenclatura "Rio Guaíba",
conforme padronização do IBGE em seus mapas oficiais \cite{ibgeguai}. Essa decisão é baseada na maior disponibilidade de
literatura relacionada à cheias fluvias e a utilização de fontes de dados provenientes do IBGE. Contudo, quando
apropriado, trabalhos envolvendo cheias em Lagos também serão utilizados de forma comparativa, como o de Güldal et al.
(2009) \cite{gudal}.

A Região Hidrográfica do Guaíba, da qual a Bacia Hidrográfica do Lago Guaíba faz parte, se extende pela parte
centro-leste do estado e inclui as Bacias Hidrográficas do Lago Guaíba, Sinos, Gravataí, Caí, Baixo-Jacuí, Alto-Jacuí,
Vacacaí, Pardo e Taquari-Antas, vide \autoref{regiao_guaiba}. A mesma constitui uma das duas Regiões Hidrográficas do
estado junto com a Região Hidrográfica do Rio Uruguai. 

Apresentando, em sua maioria, um padrão de drenagem dendrítico — caracterizado por uma rede de canais que se assemelha
aos ramos de uma árvore e que segue de forma adaptada ao relevo, vide \autoref{drain}.1 —, a Região Hidrográfica do
Guaíba concentra o escoamento das águas no Rio Guaíba, configurando um sistema de convergência típica em bacias com
vales em “V” e substrato predominantemente impermeável. Embora o padrão dendrítico seja predominante, variações locais
ocorrem devido à extensão e diversidade geomorfológica da região. A deposição de sedimentos, especialmente durante os
períodos de cheia, é marcante.

\begin{figure}[htb]
	\caption{\label{regiao_guaiba} Mapa da Região Hidrográfica do Guaíba}
	\begin{center}
	    \includegraphics[scale=0.35]{img/regiao_guaiba.png}
	\end{center}
	\legend{Fonte: Secretaria do Meio Ambiente e Infraestrutura do Rio Grande do Sul \cite{sema}}
\end{figure}

\subsection{Enchentes históricas}

A região da Bacia Hidrográfica do Guaíba é historicamente suscetível à cheias. A mais notável destas foi a Enchente de
1941 - marco hidrológico que permanece como referência catastrófica. Entre de Abril e Maio daquele ano, um período de
alta precipitação atingiu o estado e, somente na capital, a precipitação chegou a 791 milímetros. Combinado com o
escoamento das águas das chuvas ao redor do estado na Bacia Hidrográfica do Guaíba, o nível do mesmo chegou a 4,76
metros acima da cota normal \cite{dep}, submergindo bairros ribeirinhos como Centro Histórico, Menino Deus e Azenha por
22 dias consecutivos. Dos 272 mil habitantes então recenseados, 70 mil (25,7\%) foram desalojados, com transporte urbano
reduzido a embarcações de emergência. A contaminação do sistema de saneamento pela água pluvial desencadeou epidemias de
doenças hídricas nos meses subsequentes, enquanto empresas levaram meses para retomar operações. Este evento foi um dos
primeiros com bons registros na região. Subsequentemente, as cheias foram mais recorrentes: em 1967 o nível atingiu 3,13
m sem proteção do Muro da Mauá (então em construção); em 1984 registrou-se 2,50 m com ativação das comportas do sistema;
e em 2002 o pico de 2,46 m aproximou-se do limite operacional da infraestrutura.

A suscetibilidade regional é estrutural. Porto Alegre ocupa uma planície aluvial crítica com 35\% da área urbana abaixo
da cota 3 m e é localizada na confluência de quatro bacias hidrográficas (Jacuí, Gravataí, Sinos, Caí) que drenam 30\%
do território gaúcho \cite{dep_topografia}, além de ser cercada por arroios e pela Lagoa dos Patos ao sul. Este gargalo
hidrodinâmico coincide precisamente com o núcleo metropolitano, dado que na estreita região entre a Usina do Gasômetro e
a Ilha da Pintada, se encontra o maior ponto de represamento da região graças ao encontro dos rios. O crescimento urbano
e o desmatamento de encostas aumentaram a impermebilidade do solo e consequentemente diminuiram a capacidade de
inflitração urbana da região, aumentando o risco de uma área já vulnerável. Isto, somado à densidade populacional
presente nestas regiões, tornam crítico o monitoramento da cota das águas do Guaíba.

Nas últimas décadas, padrões pluviais intensificados pelas mudanças climáticas geraram eventos cada vez mais extremos:
em 2023 precipitações geradas por um ciclone extratropical levaram a dois eventos extremos no estado, um em Setembro e o
outro em Novembro. Afetando principalmente a Região do Vale do Taquari, na qual atingiu 359 mil pessoas, o excesso de
água causado pelas chuvas também chegou ao Rio Guaíba. No dia 21 de Novembro, a cota do rio chegou á 3,46 metros, o mais
alto desde a enchente de 1941. Porém, neste caso, os sistemas de prevenção de cheias da cidade consguiram parar grande
parte do avanço das águas.

\subsection{Sistemas de prevenção de cheias na RMPOA}
\lipsum[1-8]
% Casas de bombas, muro da Mauá, Galerias pluviais e como a falta de manutenção fez esses sistemas falharem -> Falta de
% Resiliência na malha urbana;

% Muro da Mauá. A enchente histórica trouxe à tona a ideia de construir um muro de proteção para Porto Alegre. Cerca de
% 30 anos depois do evento, entre 1971 e 1974, foi construído o Muro da Mauá, uma estrutura de concreto de 3 metros de
% altura e 2,6 km de extensão que faz parte de um sistema de contenção de águas.

% Diques. Atualmente, a cidade tem 24 km de diques que margeiam a FreeWay até a Avenida Castelo Branco e seguem pela
% Avenida Beira Rio; e mais 44 km de diques internos, às margens dos arroios que atravessam a cidade. São 68 km de
% diques, 14 comportas e 19 casas de bomba, de acordo com a prefeitura de Porto Alegre

\subsection{Enchentes de 2024}
\lipsum[1-6]

% INFORMATIVOS DO INMET -> https://portal.inmet.gov.br/informativos# (Tem que pesquisar no google cada um deles);

% Vale mensionar as enchentes de 2023 como introdução -> Talvez previsões de aumentos de enchentes na região sudeste do
% país? Tem um estudo que foi realizado no governo da Dilma;

% https://www.ncei.noaa.gov/monitoring-content/sotc/global/2024/apr/Brazil-Floods-RioGrandeSul-202404.pdf
% https://web.archive.org/web/20240930134441/https://www.nytimes.com/2024/05/02/world/americas/brazil-rain-floods.html
% https://www.bbc.com/news/world-latin-america-68935249
% https://www.reuters.com/graphics/BRAZIL-RAINS/FLOODS/egpbazlzbvq/

% Cronologia: https://www.bbc.com/portuguese/articles/cd1qwpg3z77o

% As enchentes em Porto Alegre exibiram características tanto de enchentes pluviais quanto fluviais, mas foi o último
% que gerou a maior parte da crise -> Foco no Guaíba. Isso se dever por ter tido a característica de uma enchente
% fluvial, isto é, o Rio foi aquilo que transbordou, mas também pode ser caracterizar como uma enchente pluvial urbana,
% no sentido de que o escoamento do Rio que levou ao aumento do Nível da Água e a falha do escoamento urbano que levou a
% água para dentro da cidade. 

% 5. Climate change, El Niño and infrastructure failures behind massive floods in southern Brazil

\chapter{Revisão bibliográfica}
Neste capítulo exploraremos os modelos escolhidos na parte experimental do projeto. Um dos maiores critérios de escolha
foi baseado em ter pelo menos um representante das principais áreas de predição de séries temporais baseadas em dados
para obter-se uma comparação entres as diferentes famílias de métodos. Logo, se viu necessária a presença de um modelo
estatístico clássico como base de comparação, modelos de Redes Neurais devido ao seus grandes usos e efetividades e, por
último, pelo menos um modelo baseado em árvores de decisão devido a sua intuitiva compreensão e permissividade em
círculos de ciência de dados. 
    \section{Análise Comparativa da Literatura}
    % Uma alternativa que visa superar essas limitações são os modelos baseados em dados, que utilizam técnicas para
% identificar e explorar relações estatísticas entre os dados de entrada e saída. Com a aprendizagem dos padrões
% presentes nos dados históricos, é possível mapear as dinâmicas subjacentes do sistema e utilizá-las para prever
% valores futuros. Assim, esses modelos são capazes de realizar previsões precisas com base em séries temporais,
% capturando tendências e comportamentos complexos que podem ser difíceis de modelar diretamente por meio de abordagens
% tradicionais. Diferentemente das mesmas, que demandam tempo e dados específicos, os modelos baseados em dados oferecem
% maior agilidade para previsões em tempo real, sendo uma escolha mais adequada para sistemas de alerta precoce e
% mitigação de desastres de cheias \cite{ta}.

% Uma das maiores vantagens de utilizar essa abordagem é o grande desenvolvimento do campo de Análise e Predição de
% Séries Temporais. Com isso, é possível utilizar as técnicas desenvolvidas para a aplicação em problemas com uma
% modelagem semelhante, mas em contextos distintos \cite{hyndman}, como Economia, Biologia, entre outros. Dentro deste
% contexto, uma ampla gama de modelagens são utilizadas para a predição de séries temporais. Isso apresenta uma
% oportunidade de análise e comparação entre as vantagens e desvantagens das mesmas com o intuíto de descobrir a mais
% efetiva para a situação particular \cite{makridakis, morales}.

% Uma abordagem mais clássica e altamente implementada destes modelos vem na forma de modelos estatísticos, como o ARIMA
% (\textit{Autoregressive Integrated Moving Average}) \cite{box}, um modelo que descreve processos estocásticos usando
% uma combinação de autoregressão, isto é, a predição de uma Série Temporal usando somente valores anteriores da mesma,
% a média móvel - ajuste das previsões considerando os erros residuais das observações anteriores, o que permite
% capturar flutuações - e, por último, a diferenciação dos valores consecutivos da série, reduzindo o impacto de padrões
% de longo prazo e tendências. Além disso, um elemento de sazionalidade, ou seja, padrões repetidos em com uma certa
% frequência fixa, como padrões mensais ou trimestrais, podem ser compreendidos através de uma extensão do ARIMA, o
% SARIMA. Historicamente, esses modelos têm sido usados em predições em contextos semelhantes como em Atashi et al.
% (2022) \cite{atashi} que comparou o desempenho do SARIMA com os modelos LSTM e \textit{Random Forest} para predição do
% nível da água no Rio Vermelho do Norte, localizado na fronteira entre os Estados Unidos e o Canadá, e Güldal et al.
% (2009) \cite{gudal} que utilizou o ARMA para a predição do Lago Eğirdir na Túrquia, oferecendo uma base para
% comparações com métodos modernos.

% Esses modelos, eficazes e de baixo custo computacional, apresentam tanto limitações quanto vantagens quando comparados
% com modelos mais recentes desenvolvidos nas áreas de Aprendizado de Máquina (\textit{Machine Learning}) e Aprendizado
% Profundo (\textit{Deep Learning}) \cite{kontopoulou}. No âmbito de \textit{Machine Learning}, modelos baseados em
% \textit{Gradient Boosting} têm se mostrado promissores em situações de desastres, onde a rapidez e precisão nas
% previsões são essenciais para mitigar impactos \cite{albahri}. Um exemplo amplamente utilizado é o XGBoost
% (\textit{eXtreme Gradient Boosting}) \cite{chen_xgboost}. O mesmo utiliza árvores de decisão em série, ajustando
% constantemente os erros das previsões anteriores. Além disso, ele incorpora regularização para reduzir o risco de
% sobreajuste (\textit{overfitting}) e otimizações específicas para lidar com grandes volumes de dados. Uma
% implementação semelhante, visando reduzir o tempo computacional, é dada pelo LightGBM \cite{lightgbm}. Ela mantém as
% vantagens do XGBoost, como otimização para dados esparsos, treinamento paralelo, múltiplas funções de perda e
% regularização, mas altera o método de construção das árvores de decisão, mostrando-se uma boa alternativa para obter
% resultados semelhantes de forma mais eficaz, algo imprescindível no contexto de uma crise climática como a ocorrida.

% Por outro lado, em Aprendizado Profundo, modelos baseados em Redes Neurais Recorrentes (RNNs), como o \textit{Long
% Short Term Memory} (LSTM), são amplamente utilizados em situações semelhantes, como demonstrado por Borwarnginn et al.
% \cite{borwarnginn}. O LSTM foi projetado especificamente para lidar com dependências de longo prazo em séries
% temporais, garantindo que informações relevantes de eventos passados sejam preservadas, ao mesmo tempo que mantém o
% contexto de eventos mais recentes. Essa análise é particularmente útil em problemas onde padrões temporais complexos
% e, geralmente, não lineares precisam ser modelados.
    \section{Seasonal Autoregressive Integrated Moving Average (SARIMA)}
    O modelo de Média Móvel Integrada Auto-Regressiva com Sazonalidade (SARIMA) é uma extensão do modelo de Média Móvel
Integrada Auto-Regressiva (ARIMA) que combina três elementos específicos em sua composição, a autoregressão (AR), a
média móvel (MA) e a diferenciação dos valores consecutivos da série (I), sendoexpandido para lidar com padrões sazonais
\cite{box}.

A componente de autoregressão descreve a dependência de uma observação atual e observações anteriores, em defasagens
$p$. Em um modelo puramente autorregressivo, uma nova observação $X_t$ é expressa como uma combinação linear dos $p$
valores anteriores, ponderados por coeficientes $\alpha_i$, conforme a equação \autoref{ar}, onde $\epsilon_t$
representa um termo de erro aleatório:

\begin{equation}
    \label{ar}
    X_t=\alpha_1X_{t-1} + \alpha_2X_{t-2} +(...) + \alpha_pX_{t-p} + \epsilon_t
\end{equation}

Já a componente de média model indica a relação de dependência de uma observação e o erro residual em pontos anteriores.
Isto é, $X_t$ depende de perturbações aleatórias anteriores $\epsilon_{t-i}$ moduladas por coeficientes $\theta_i$,
conforme \autoref{ma}:

\begin{equation}
    \label{ma}
    X_t=\theta_1\epsilon_{t-1} + \theta_2\epsilon_{t-2} + (...) + \theta_q\epsilon_{t-q}
\end{equation}

A combinação dos dois componentes forma o modelo ARMA, como mostrado na equação \autoref{arma}:

\begin{equation}
    \label{arma}
    X_t = \epsilon_t + \sum_{i=1}^{p}\phi_i X_{t-i} + \sum_{i=1}^q \theta_i \epsilon_{t-i}
\end{equation}

Entretanto, modelos ARMA pressupõem que a série seja estacionária (com média e variância constantes ao longo do
tempo). Na prática, muitas séries reais apresentam tendências. Para tratar essa não-estacionariedade, o modelo ARIMA
introduz o parâmetro $d$, referente à ordem de integração.

O parâmetro $d$ indica o número de vezes que a série deve ser diferenciada para tornar-se estacionária. A primeira
diferença ($d=1$) é definida como $\Delta X_t = X_t - X_{t-1}$. Para facilitar a manipulação matemática, utiliza-se o
operador de defasagem (ou \textit{lag}) $L$, definido como $L^i X_t = X_{t-i}$. Assim, a diferenciação de ordem $d$ é
representada pelo operador $(1 - L)^d$.

Ao aplicar este operador à série original, transformamos $X_t$ em uma série estacionária antes de aplicar os componentes
AR e MA. O modelo ARIMA($p, d, q$) resultante é expresso na \autoref{arima}:

\begin{equation}
    \label{arima}
    \left(1 - \sum_{i=1}^p \phi_i L^i\right) (1 - L)^d X_t = \left(1 + \sum_{i=1}^q \theta_i L^i\right) \epsilon_t
\end{equation}

A extensão SARIMA incorpora componentes sazonais adicionais para lidar com padrões que se repetem em intervalos fixos,
como meses ou trimestres. Esses termos sazonais seguem a mesma estrutura do ARIMA, com parâmetros $P$, $D$, $Q$ e
período de sazonalidade $s$. Essa estrutura permite ao modelo capturar tanto as dependências de curto prazo quanto os
padrões cíclicos de longo prazo presentes em séries temporais com sazonalidade definida.

% \begin{equation} \label{sarima} \Phi_P(L^s) \phi_p(L)(1 - L)^d (1 - L^s)^D X_t = \Theta_Q(L^s) \theta_q(L) \epsilon_t
%     \end{equation}

% Nessa expressão $\phi_p(L)$ e $\theta_q(L)$ são os polinômios não sazonais AR e MA, respectivamente, $\Phi_P(L^s)$ e
% $\Theta_Q(L^s)$ são os polinômios sazonais AR e MA, aplicados com defasagem múltipla de $s$ períodos e $(1 - L)^d$
% representa a diferenciação não sazonal e $(1 - L^s)^D$ a diferenciação sazonal.
    \section{Long Short Term Memory (LSTM)}
    O Memória de Curto e Longo Prazo (Long Short-Term Memory — LSTM) é uma variante das Redes Neurais Recorrentes (Recurrent
Neural Networks — RNN) projetada para resolver o problema do gradiente desaparecendo (vanishing gradient problem) que
afeta as RNNs quando precisam aprender dependências de longo prazo em sequências extensas. Enquanto uma RNN padrão
transmite apenas seu estado oculto (hidden state) de um passo ao outro, sofrendo perda de informações importantes ao
longo do tempo, o LSTM introduz um estado interno ou célula de memória que pode reter informações relevantes por vários
\textit{timesteps}, de acordo com sua importância.

Cada célula de LSTM é composta por quatro camadas neurais distintas: três portões — portão de esquecimento (forget
gate), portão de entrada (input gate) e portão de saída (output gate) — e uma camada responsável por gerenciar o próprio
estado interno. Esses portões combinam funções sigmoide e multiplicações ponto a ponto para controlar de forma
diferenciada o que deve ser descartado, adicionado ou propagado ao passo seguinte. A figura a seguir compara o fluxo de
informação em uma RNN simples e em uma célula de LSTM.  
% FIGURE: insira “LSTM3-SimpleRNN.png” e “LSTM3-chain.png” para comparação RNN × LSTM

O portão de esquecimento (forget gate) decide quais componentes de \(C_{t-1}\), o estado interno anterior, serão
removidos:
\[
f_t = \sigma\bigl(W_f[h_{t-1}, x_t] + b_f\bigr),
\]
onde \(\sigma\) é a sigmoide, \(h_{t-1}\) o estado oculto anterior, \(x_t\) a entrada atual, e \(W_f, b_f\) pesos e
vieses. Valores de \(f_t\) próximos a 0 indicam descarte, e próximos a 1 indicam retenção.  
% FIGURE: insira “LSTM3-focus-f.png” ilustrando o portão de esquecimento

O portão de entrada (input gate) regula a incorporação de novas informações através de duas etapas:  
1. Um cálculo sigmoide determina as posições a serem atualizadas:
   \[
   i_t = \sigma\bigl(W_i[h_{t-1}, x_t] + b_i\bigr).
   \]
2. Uma camada \(\tanh\) gera o vetor de candidatos:
   \[
   \widetilde{C}_t = \tanh\bigl(W_c[h_{t-1}, x_t] + b_c\bigr),
   \]
   com valores em \([-1,1]\).  
% FIGURE: insira “LSTM3-focus-i.png” ilustrando o portão de entrada

Em seguida, o estado interno é atualizado por
\[
C_t = f_t \,\odot\, C_{t-1} \;+\; i_t \,\odot\, \widetilde{C}_t,
\]
onde \(\odot\) denota multiplicação ponto a ponto.  
% FIGURE: insira “LSTM3-focus-C.png” mostrando a atualização de \(C_t\)

Finalmente, o portão de saída (output gate) define a porção de \(C_t\) que se tornará o próximo estado oculto:
\[
o_t = \sigma\bigl(W_o[h_{t-1}, x_t] + b_o\bigr),
\qquad
h_t = o_t \,\odot\, \tanh(C_t).
\]
% FIGURE: insira “LSTM3-focus-o.png” ilustrando o portão de saída

As equações essenciais de um LSTM resumem-se a:
\[
\begin{aligned}
f_t &= \sigma(W_f[h_{t-1}, x_t] + b_f),\\
i_t &= \sigma(W_i[h_{t-1}, x_t] + b_i),\\
\widetilde{C}_t &= \tanh(W_c[h_{t-1}, x_t] + b_c),\\
C_t &= f_t \,\odot\, C_{t-1} + i_t \,\odot\, \widetilde{C}_t,\\
o_t &= \sigma(W_o[h_{t-1}, x_t] + b_o),\\
h_t &= o_t \,\odot\, \tanh(C_t).
\end{aligned}
\]

A função de ativação \(\tanh\) pode ser substituída por outras, como ReLU, conforme a aplicação. Graças à combinação de
portões e do estado interno, o LSTM retém informações críticas por longos intervalos e descarta dados irrelevantes,
superando com eficiência o problema do gradiente desaparecendo e viabilizando o aprendizado de dependências de longo
prazo em tarefas de séries temporais e processamento de linguagem natural.



% % \section*{Referências} \begin{itemize} \item \url{https://colah.github.io/posts/2015-08-Understanding-LSTMs/} \item
% % \url{https://d2l.ai/chapter_recurrent-modern/lstm.html} \item %
% \url{https://www.tensorflow.org/guide/keras/sequential_model} \end{itemize}

    \section{Gradient Boosting Machine}
    O \textit{Gradient Boosting} é um método de aprendizado de máquina baseado em conjuntos de árvores de decisão que
constrói modelos preditivos sequencialmente. Esta técnica combina modelos fracos (tipicamente árvores rasas) para formar
um preditor robusto, onde cada novo modelo corrige os erros residuais do conjunto anterior \cite{friedman}. O processo
inicia-se com um modelo constante $F_0(x) = min_{\gamma} \sum_{i=1}^n \mathcal{L}(y_i, \gamma)$, geralmente a média
dos valores-alvo. 

Em cada iteração $m$, o algoritmo calcula o gradiente negativo da função de perda $\mathcal{L}$ em relação às predições
atuais: $r_{im} = -\left[ \frac{\partial \mathcal{L}(y_i, F_{m-1}(x_i))}{\partial F_{m-1}(x_i)} \right]$.
Subsequentemente, ajusta-se uma árvore de decisão $h_m(x)$ a esses resíduos e atualiza o modelo conforme $F_m(x) =
F_{m-1}(x) + \gamma_m h_m(x)$, onde $\gamma_m$ denota a taxa de aprendizado. 

Entre suas vantagens destacam-se a alta precisão preditiva, flexibilidade para dados heterogêneos, suporte a funções de
perda personalizadas e aplicabilidade em regressão e classificação. Contudo, apresenta sensibilidade a
\textit{overfitting} (controlável via \textit{early stopping}), demanda ajuste rigoroso de hiperparâmetros como
profundidade das árvores, taxa de aprendizado e número de iteradores, além de custo computacional elevado.
    \section{Revisão das Métricas}
    \subsection{Root Mean Squared Error}
É um estimador da diferença entre o valor observado e o valor real; Dado pela seguinte fórmula $$ RMSE =
\sqrt{\frac{1}{n}\sum_{i=1}^n(Y_i - Y_i')^2}$$ Yi são os valores reais e Y'i são os valores observados. Além disso, *n*
é o número de amostras. Aquilo que é medido na verdade é a média do erro ao quadrado, ou seja, outliers são punidos mais
fortemente do que outros pontos, pois quanto mais longe estiver, maior será o RMSE;


\subsection{Mean Absolute Error}


\subsection{Nash-Sutcliffe Model Efficiency Coefficient}
Nash Sutcliffe efficiency -> é uma Machine Learning Metrics|Métrica que compara a magnitude residual do resíduo em
comparação com a variação dos dados medidos. Ou seja, ele server para compara dados simulados (modelo) com dado medido
(real). É dado por (Com $Q_0$ como média dos dados observados e $Q_m$ os dados modelados e $Q^t_0$ os dados observados
em um dado tempo):

\[
\begin{aligned}
NSE = 1 - \frac{\sum^T_{t=1} (Q^t_0 - Q^t_m)^2}{\sum^T_{t=1} (Q^t_0 - Q_0)^2}
\end{aligned}
\]

Essa métrica é normalizada entre 0 e 1. Para o caso de 1, o modelo perfeitamente descreve os dados observados. Contudo,
o que é interessante dessa métrica é para modelos com NSE 0. Para esses casos, o modelo teve a mesma habilidade
preditiva do que a média de uma Série Temporal em relação à soma do erro quadrado; * Valores mais próximos de 1 sugerem
uma habilidade preditiva maior; Como sempre, essa normalização não é perfeita. Para casos com NSE<0, a média do
observado é um melhor preditor do que o modelo;

% \part{Preparação da pesquisa}
\chapter{Planejamento da pesquisa}
Neste capítulo, será descrito de forma prática como que a pesquisa foi realizada e quais ferramentas foram utilizadas
no processo de análise comparativa entre os modelos preditivos. 
\section{Metodologia}
Este trabalho implementou uma abordagem experimental para avaliação comparativa dos modelos preditivos descritos
anteriormente. Inicialmente, os dados do SNIRH foram submetidos a um pré-processamento abrangente, incluindo tratamento de
valores faltantes mediante técnicas de interpolação temporal, detecção de \textit{outliers} via análise por PCA e ECOD, e
normalização. Os dados foram particionados cronologicamente em conjuntos de treino (70\%), validação (15\%) e
teste (15\%), assegurando-se a preservação da sequência temporal para prevenir vazamento de informação.

Cada categoria de modelo exigiu transformações específicas dos dados, visto que cada um conta com um algoritmo
diferente de implementação e uma lógica interna distinta. Para o SARIMA, foi aplicada diferenciação e decomposição
sazonal. Modelos baseados em árvores requeriram codificação de variáveis temporais. Já redes LSTM demandaram janelamento
sequencial. O treinamento foi conduzido em \textit{pipelines} isolados, incorporando validação cruzada temporal e
otimização Bayesiana de hiperparâmetros para minimizar o KGE no conjunto de validação.

Na fase de avaliação, as predições geradas para o período de teste foram confrontadas com dados observados reservados de
teste mediante métricas hidrologicamente relevantes - coeficiente de Nash-Sutcliffe (NSE), RMSE e erro absoluto médio
(MAE) - complementadas por análise de sensibilidade a eventos extremos e avaliação de robustez computacional. Esta
estrutura metodológica visou garantir comparação objetiva do desempenho preditivo.

% Os resultados da análise comparativa serão disponibilizados através de um painel interativo (dashboard) desenvolvido
% após a consolidação das avaliações dos modelos. Esta plataforma de visualização, hospedada em acesso aberto, integrará
% dados históricos do nível do Rio Guaíba com as predições geradas pelos diferentes modelos, apresentando múltiplos
% cenários hidrológicos contextualizados por análises explicativas. Os usuários poderão explorar dinamicamente os dados
% mediante filtros temporais personalizáveis, permitindo focalizar períodos de interesse específico como eventos extremos
% ou fases sazonais críticas.

\section{Fonte de Dados}
Este trabalho empregou dados primários do Portal Hidroweb, integrante do Sistema Nacional de Informações sobre Recursos
Hídricos (SNIRH) da Agência Nacional de Águas (ANA) \cite{snirh}. Este sistema consolida informações de estações
hidrometeorológicas, incluindo níveis fluviais, vazões, precipitação, entre outros. A coleta automatizada dos dados foi
viabilizada mediante acesso à API do Hidroweb, concedido pela ANA ao autor.

As possíveis estações foram obtidas do endpoint \textit{HidroInventarioEstacoes} \cite{HidroInventarioEstacoes},
enquanto as séries detalhadas foram obtidas do endpoint \textit{HidroinfoanaSerieTelemetricaDetalhada}
\cite{HidroinfoanaSerieTelemetricaDetalhada}.

As estações analisadas pertencem à Bacia Hidrográfica do Guaíba, no Rio Grande do Sul, abrangendo em especial o Rio
Guaíba ('87200000') e os principais rios afluentes — Taquari ('86001000'), Caí ('87220000'), Gravataí ('87240000'),
Sinos ('87230000') e Jacuí ('85001000')— a fim de garantir ampla cobertura espacial e boa qualidade dos dados. Para
isto, foi necessário investigar todas as possíveis estações para estes rios. O total inicial era de 258 estações. Após
filtrar por estações em operação e estações telemétricas, restaram apenas 43 estações. As séries históricas destas foram
analisadas quanto à disponibilidade, qualidade, lacunas, duração e representatividade. Também excluíram-se estações com
finalidades específicas (energéticas, piezométricas etc.), que não apresentam as informações relevantes para este
estudo. 

\begin{figure}[htb] \caption{\label{funil} Funil de estações selecionadas do SNIRH} \begin{center}
  \includegraphics[scale=0.34]{img/funil_estações.png} \end{center} \legend{Fonte: Reprodução própria} \end{figure} 

Com isto, determinou-se que somente 15 estações na Bacia Hidrográfica do Guaíba apresentavam informações da forma
desejada. Destas, foram selecionadas pelo menos uma para cada um dos principais rios listados acima. Alguns rios, como o
Sinos e o Taquari, apresentavam 3 estações cujas informações poderiam ser utilizadas, mas por fim, foi determinado que
no máximo 2 estações seriam selecionadas de cada rio, a fim de manter um equilíbrio de dados e não introduzir um maior
peso para estes corpos hídricos, inclusive quando alguns, como o Rio Jacuí, só teriam 1 estação disponível.

Estas foram coletadas com dados à cada 15 minutos para um período histórico de 5 anos. Para o próprio Guaíba, adotou-se
como alvo, a estação CAIS MAUÁ C6 ('87450004') e complementarmente a estação do USINA DO GASÔMETRO - CHEIA 2024
('87444000'), estratégia necessária devido à descontinuidade nos registros da primeira estação em 2-3 de maio de 2024.
Conforme atestado pela SEMA \cite{informativo_sema}, este hiato decorreu de danos nos equipamentos de medição durante as
enchentes.

Por fim, a lista de estações selecionadas foram:
\begin{table}[h!]
    \centering
    \caption{Estações Hidrometeorológicas por Bacia}
    \label{tab:estacoes_hidro}
    \begin{tabular}{|l|l|l|}
        \hline
        \textbf{Rio} & \textbf{Estação} & \textbf{Código Estação} \\
        \hline
        \hline
        \multirow{2}{*}{\textbf{Guaíba}} & CAIS MAUÁ C6 & 87450004 \\
        \cline{2-3}
        & USINA DO GASÔMETRO (CHEIA 2024) & 87444000 \\
        \hline
        \multirow{2}{*}{\textbf{Taquari}} & MUÇUM & 86510000 \\
        \cline{2-3}
        & ENCANTADO & 86720000 \\
        \hline
        \multirow{2}{*}{\textbf{Caí}} & LINHA GONZAGA & 87150000 \\
        \cline{2-3}
        & BARCA DO CAÍ & 87170000 \\
        \hline
        \textbf{Gravataí} & PASSO DAS CANOAS - AUXILIAR & 87399000 \\
        \hline
        \multirow{2}{*}{\textbf{Sinos}} & SÃO LEOPOLDO & 87382000 \\
        \cline{2-3}
        & CAMPO BOM & 87380000 \\
        \hline
        \textbf{Jacuí} & RIO PARDO & 85900000 \\
        \hline
    \end{tabular}
\end{table}

Para uniformizar as séries, adotou-se `2018-08-01` como data inicial comum e `2025-12-01` como data final comum,
totalizando cerca de seis anos de dados telemétricos, totalizando mais de 2 milhões de registros. Devido ao limite de 30
dias imposto pela API, o processo de coleta foi realizado em chunks mensais posteriormente concatenados.

Foram coletados todos os dados disponíveis para estas estações neste período, incluindo dados do nível da água (e.g.,
\textit{Cota\_Adotada}), dados meteorológicos (e.g., \textit{Chuva\_Adotada}), dados descritivos da estação (e.g.,
\textit{Altitude}) e indicadores de controle de qualidade de cada variável, identificados pelo sufixo \textit{\_Status}.
Estas variáveis de status acompanham as variáveis de medição e seguem um padrão determinado pelo SNIRH que foi expandido
para utilizar estas variáveis como indicadores não só da qualidade dos dados, mas também como indicadores de variáveis
faltantes ou preenchidas. Este padrão é dado de tal forma: Status 0=Normal, 1=Suspeito, 2=Ruim, 3=Muito ruim,
4=Preenchido/Faltante, 5=Valor Atípico (outliers);

% Bateria, Chuva_Acumulada, Chuva_Acumulada_Status, Chuva_Adotada, Chuva_Adotada_Status, Cota_Adotada,
% Cota_Adotada_Status, Cota_Display, Cota_Display_Status, Cota_Manual, Cota_Manual_Status, Cota_Sensor,
% Cota_Sensor_Status, Data_Atualizacao, Data_Hora_Medicao, Pressao_Atmosferica, Pressao_Atmosferica_Status,
% Temperatura_Agua, Temperatura_Agua_Status, Temperatura_Interna, Vazao_Adotada, Vazao_Adotada_Status, codigoestacao,
% Temperatura_Interna_Status

\section{Processamento de dados}
O primeiro e maior desafio enfrentado nos dados coletados foram os valores faltantes. Certos atributos estão
consistentemente faltando do dataset, como valores de 'Pressao\_Atmosferica' e 'Temperatura\_Agua' infelizmente
apresentam muitos valores faltantes para serem de qualquer uso nessa análise e logo serão desconsiderados daqui pra
frente.

\begin{figure}[htb] \caption{\label{faltantes} Valores faltantes} \begin{center}
  \includegraphics[scale=0.53]{img/faltantes.png} \end{center} \legend{Fonte: Reprodução própria} \end{figure} 

Também vale considerar todos aqueles relacionados à Cota. Estes, em especial, são formas de medir o nível da água.
'Cota\_Manual', por exemplo, representa medidas feitas com uma régua ao invés das medidas automatizadas do sensor. Dado
isto, usaremos valores de 'Cota\_Manual' e 'Cota\_Sensor' para preencher valores faltantes de 'Cota\_Adotada';

Para escolher quais dados seriam utilizados nos modelos de predição, também foi realizada uma análise da correlação dos
valores, presente na \autoref{correlacao}.

\begin{figure}[htb] \caption{\label{correlacao} Correlação entre os dados das estações} \begin{center}
  \includegraphics[scale=0.275]{img/correlacao.png} \end{center} \legend{Fonte: Reprodução própria} \end{figure} 

A análise de correlação evidencia que as variáveis 'Cota\_Manual', 'Cota\_Sensor' e 'Cota\_Display' apresentam forte
correlação linear com o alvo 'Cota\_Adotada'. Observa-se que 'Cota\_Manual' possui mais de 99\% de valores ausentes, e
os registros válidos restantes concentram-se majoritariamente em pontos nos quais 'Cota\_Adotada' não está disponível.
Esse padrão sugere que as medições manuais são utilizadas como mecanismo de contingência em situações de manutenção ou
falha dos sensores automáticos. Adicionalmente, 'Cota\_Sensor' e 'Cota\_Display' mostram-se, na maior parte do tempo,
virtualmente idênticas a 'Cota\_Adotada', o que é consistente com o fluxo operacional da estação, no qual as leituras
oficiais derivam predominantemente dos sensores automáticos.

A variável 'Vazao\_Adotada' apresenta elevada correlação com 'Cota\_Adotada', o que é consistente com o comportamento
hidrodinâmico esperado: incrementos de vazão (m\textsuperscript{3}/s) em um canal de geometria aproximadamente fixa
tendem a elevar o nível d'água até o ponto em que o escoamento interage áreas de várzea adjacentes. Essa variável
constitui um preditor importante devido à sua forte associação com a variável alvo e à sua baixa correlação com o
restante dos atributos, indicando aporte de informação complementar ao modelo. No entanto, há uma alta taxa de valores
ausentes (>15\%), cuja magnitude varia entre estações. A estação alvo (Guaíba) não possui registros de vazão, o que
inviabiliza seu uso direto nessa localidade. Ainda assim, dada a elevada dependência física entre vazão e nível,
optou-se por manter essa variável no processo de modelagem, mesmo que as mesmas impliquem um maior processamento dos
dados.

Com base na análise de correlação e na avaliação de valores faltantes, foram selecionadas as seguintes variáveis para os
modelos de \textit{Machine Learning}: 'Cota\_Adotada', correspondente ao nível do rio em centímetros, cujos valores
ausentes serão imputados a partir de 'Cota\_Manual' e 'Cota\_Sensor' sempre que disponíveis; 'Chuva\_Adotada',
representando a precipitação acumulada desde a última leitura, em milímetros; 'Temperatura\_Interna', Correspondente à
temperatura atmosférica medida no equipamento, em °C; e 'Vazao\_Adotada', que indica o fluxo de água registrado pelo
sensor em $\frac{m^{3}}{s}$.

A variável 'Chuva\_Acumulada' também foi inicialmente escolhida e trabalhada junto das outras variáveis por um bom
tempo. Contudo, depois de análises mais aprofundadas, a parte "Acumulada" da mesma se mostrou inconsistênte, dado que
deveria representar a precipitação acumulada na estação por mês, mas o mesmo não se mostrou consistênte através das
diferentes estações. Logo, junto de outras transformações, a mesma foi recriada a partir da variável 'Chuva\_Adotada',
em um processo que será explicado melhor mais a frente.

\section{Ferramentas e recursos utilizados}
O desenvolvimento da parte escrita deste trabalho será feita em \LaTeX, utilizando o modelo \abnTeX, desenvolvido por
Araujo et al (2015) \cite{abntex2classe} de acordo com as regras da ABNT NBR 14724:2011, ABNT NBR 6024:2012 e outras.

A implementação dos modelos será desenvolvida em Python, selecionado por seu ecossistema robusto para análise de dados e
aprendizado de máquina. A biblioteca \textit{statsmodels} será empregada para modelagem SARIMA, enquanto redes LSTM
serão construídas utilizando a API de baixo nível do \textit{TensorFlow}. Para os modelos baseados em árvores, serão
adotadas as bibliotecas nativas do XGBoost e LightGBM. Uma arquitetura modular será implementada mediante classes
especializadas para cada modelo, integrando \textit{pipelines} do \textit{scikit-learn} para garantir fluxos
reprodutíveis de pré-processamento. O processamento numérico e manipulação de séries temporais serão realizados com
\textit{NumPy} e \textit{Pandas}, respectivamente, assegurando eficiência computacional e tratamento dos dados.

O painel interativo utilizará \textit{Streamlit} para conversão ágil de scripts Python em aplicações web visuais,
garantindo integração nativa com bibliotecas científicas (\textit{Pandas}, \textit{Matplotlib}, \textit{Plotly}). A
hospedagem no \textit{Streamlit Community Cloud} assegurará disponibilidade pública contínua e atualizações automáticas
conforme novos dados ou modelos forem incorporados.

\section{Setup Experimental para treinamento de modelos} 
The training process was conducted in a machine with the following specifica- tions: 12th Gen Intel(R) Core(TM)
i5-12500H, 16GB RAM and a NVIDIA GeForce RTX 3050 Laptop 4GB VRAM. The LSTM model takes in training around 1.5 seconds
per epoch, which led to around 2.5 to 4 minutes per experiment. Additionally, the feature importance calculation for
each prediction takes around 1 second per prediction.
% \part{Resultados}

\chapter{Resultados}
\section{Comparação entre modelos}
\lipsum[1]

\section{Análise de Métricas}
\lipsum[1]

% ----------------------------------------------------------
% Finaliza a parte no bookmark do PDF
% para que se inicie o bookmark na raiz
% e adiciona espaço de parte no Sumário
% ----------------------------------------------------------
\phantompart

% ---
% Conclusão
% ---
\chapter{Conclusão}
% ---

O escopo deste trabalho revelou-se seu maior desafio. Embora a base do SNIRH seja ampla, apresentou diversos problemas
de qualidade, inconsistências e falhas nas séries temporais. O esforço necessário para corrigir e padronizar esses dados
consumiu uma parcela significativa do tempo, que poderia ter sido direcionada para etapas centrais da modelagem, como
engenharia de atributos mais profunda, testes adicionais de arquiteturas, refinamento sistemático das predições, ou a
inclusão de outras fontes de dados.

A adoção de uma arquitetura em paralelo também se mostrou complexa. Essa escolha foi motivada por um cenário realista,
no qual um sistema operacional de previsão de cheias não pode depender de um único modelo, mas deve integrar diferentes
arquiteturas e comparar múltiplos cenários, potencialmente por meio de um \textit{ensemble}. Contudo, treinar
simultaneamente modelos distintos com entradas idênticas e configurações otimizadas, excedeu as capacidades de uma
máquina comercial. Um ambiente computacional mais robusto certamente produziria resultados superiores. Mais memória RAM
permitiria ampliar o conjunto de atributos selecionados; uma GPU de maior capacidade proporcionaria paralelismo mais
efetivo e tempos de treinamento substancialmente menores. Ainda assim, por tratar-se de um estudo de caso, foi
necessário trabalhar com os recursos disponíveis.

Os resultados obtidos foram satisfatórios, já que todos os modelos superaram a baseline e apresentaram capacidade
preditiva consistente. A reprodução dos padrões gerais de variação do nível d'água também foi robusta. Embora nenhum
modelo tenha capturado adequadamente os picos extremos das enchentes, isso era esperado, visto que tal evento foi
extraordinário e não possuia análogos nas séries históricas utilizadas para o treinamento. Ainda assim, a habilidade dos
modelos em identificar aumentos e quedas repentinas mostrou-se notável.

Entre as abordagens avaliadas, os modelos baseados em \textit{Gradient Boosting} demonstraram maior potencial,
especialmente pela combinação de desempenho e velocidade de treinamento. Além disso, esses modelos oferecem margem para
aperfeiçoamentos adicionais, notadamente no tratamento do escalonamento das variáveis, cuja aplicação pode interagir de
forma indesejada com métodos baseados em árvores, uma vez que seus mecanismos internos podem ser capazes de lidar com
este problema.

% ----------------------------------------------------------
% ELEMENTOS PÓS-TEXTUAIS
% ----------------------------------------------------------
\postextual
% ----------------------------------------------------------

% ----------------------------------------------------------
% Referências bibliográficas
% ----------------------------------------------------------
\bibliography{references}

% ----------------------------------------------------------
% Glossário
% ----------------------------------------------------------
%
% Consulte o manual da classe abntex2 para orientações sobre o glossário.
%
%\glossary

% ----------------------------------------------------------
% Apêndices
% ----------------------------------------------------------

% % ---
% % Inicia os apêndices
% % ---
% \begin{apendicesenv}

% % Imprime uma página indicando o início dos apêndices
% % \partapendices

% % ----------------------------------------------------------
% \chapter{Quisque libero justo}
% % ----------------------------------------------------------

% \lipsum[50]

% \end{apendicesenv}
% % ---


% % ----------------------------------------------------------
% % Anexos
% % ----------------------------------------------------------

% % ---
% % Inicia os anexos
% % ---
% \begin{anexosenv}

% % Imprime uma página indicando o início dos anexos
% % \partanexos

% % ---
% \chapter{Morbi ultrices rutrum lorem.}
% % ---
% \lipsum[30]

% \end{anexosenv}

%---------------------------------------------------------------------
% INDICE REMISSIVO
%---------------------------------------------------------------------
\phantompart
\printindex
%---------------------------------------------------------------------

\end{document}